% !TeX root = ../../Book5.tex
% 纲要标题：## 第5章　Artin 诱导、Brauer 诱导与群论应用
% 纲要位置：工作记录/计划纲要.md，第 3845 行。
% 纲要细目（写作范围）：
% > 本章顺读：表示环与整性 → Artin 有理诱导 → Brauer 整诱导的陈述与证明 → Burnside 可解性定理 → 群论应用。Burnside 证明只使用本章整性工具和此前有限群、特征标理论，不以前面整诱导的长证明为前提。

\chapter{Artin 诱导、Brauer 诱导与群论应用}
\label{b5:ch:05}
\BookChapterNumber{b5:ch:05}

\begin{chapterabstract}
	本章研究从子群诱导特征标的有理与整数表达，证明 Artin 和 Brauer 诱导定理，并由特征标整性推得 Burnside 的 \(p^a q^b\) 阶群可解性定理。
\end{chapterabstract}

\begin{learninggoals}
	\begin{itemize}
		\item 使用表示环、次数及不可约系数，区分实际特征标与虚特征标。
		\item 证明特征标值与类和标量的整性，并说明次数整除的来源。
		\item 区分有理系数诱导、有理值特征标及仅从平凡表示诱导的不同结论。
		\item 掌握 \(p\)-初等子群的定义，逐步排除整诱导中的每个素数分母。
		\item 用 Clifford 理论证明 \(p\)-群的单项式性，将整诱导化为来自一维表示的诱导。
		\item 从素数幂大小的共轭类构造正规子群，证明 Burnside 可解性定理，并判断哪些群满足其条件。
	\end{itemize}
\end{learninggoals}

设 \(G\) 为有限群。其复正则特征标从平凡子群诱导而来；它在单位元的值是群阶，在其余元素处却全为零。若只从这个子群诱导，便看不出别的不可约特征标在非单位元处的取值。改从较大的子群诱导，所得函数会在更多共轭类上留下信息；但究竟要允许哪些子群、哪些系数，才能组合出所需的特征标？有理系数与整数系数给出不同的答案。这个问题引出 Artin 与 Brauer 两种诱导定理，也迫使我们认真检查特征标值的整性及诱导在各共轭类上的行为。

\ProseParagraphBreak
本章使用前四章的复特征标、诱导公式、张量恒等式及 Clifford 对应，也使用有限群的 Sylow 定理。所需代数整数与有限交换环的引理在章内证明。\secref{b5:sec:0504}从特征标计算转入表示环的局部化与极大理想检测，证明 Brauer 整诱导定理；\secref{b5:sec:0505}再用整性方法证明 Burnside 定理。

\ProseParagraphBreak
参考阅读中，Etingof 等的《Introduction to Representation Theory》\cite[Sections 5.2--5.4; Section 5.26]{EtingofEtAl2011RepresentationTheory}分别讨论整性、Burnside 证明和 Artin 诱导；Serre 的《Linear Representations of Finite Groups》\cite[Chapter 10, pp. 74--80]{Serre1977FiniteRepresentations}集中讨论 Brauer 诱导。本章逐素数证明使用局部化表示环 \(R(G)_{(p)}\) 的极大理想的求值描述，并把相关环论步骤展开。

% !TeX root = ../../Book5.tex
% 纲要标题：### 5.1 表示环、代数整数性与类和
% 纲要位置：工作记录/计划纲要.md，第 4169 行。
% 纲要细目（写作范围）：
% * 表示环
% * 特征标值的代数整数性
% * 诱导法处理表示环问题
% * 分圆域与特征标的域
% * 从有限阶矩阵的单位根特征值证明特征标值的代数整数性，再建立群代数中心类和及其标量作用所需的整性引理
% * 所需代数整数的和、积与有理代数整数判别在本节就地准备，不把第三卷整扩张理论作为首次学习有限群表示的前提

\section{表示环、代数整数性与类和}
\label{b5:sec:0501}

特征标的加法来自直和，乘法来自张量积。允许相减以后，诱导问题便可以放进一个环中讨论。整数系数的诱导比有理系数的诱导要求更强，需要研究特征标格、单位根和有限生成的整系数模。本节介绍这些工具。

\subsection{表示环}
\label{b5:subsec:0501:01}

\begin{definition}\label{b5:def:complex-representation-ring}
	设 \(G\) 为有限群，\(\operatorname{Irr}(G)\) 为其复不可约特征标的集合。所有有限和
	\[
		R_{\mathbb C}(G)\coloneqq \bigoplus_{\chi\in\operatorname{Irr}(G)}\mathbb Z\chi
	\]
	在类函数的逐点加法、乘法下组成的环，称为 \(G\) 的复\term[biaoshihuan]{表示环}[representation ring]，以下简记为 \(R(G)\)。其元素称为\term[xutezhengbiao]{虚特征标}[virtual character]。乘法单位是平凡特征标 \(1_G\)，而不是正则特征标。
\end{definition}
\symidx{\(\operatorname{Irr}(G)\)：不可约复特征标的集合}
\symidx{\(R_{\mathbb C}(G),R(G)\)：有限群的复表示环}

乘积确实仍在上述整数格内：若 \(V,W\) 为复表示，则 \(\chi_{V\otimes W}=\chi_V\chi_W\)，将 \(V\otimes W\) 分解成不可约表示便得到非负整数系数。每个虚特征标都能写成 \(\chi_V-\chi_W\)；此记号只记录直和的形式差，不声称存在一个维数可能为负的向量空间。

\ProseParagraphBreak
次数映射 \(\dim:R(G)\rightarrow\mathbb Z\)、\(\alpha\mapsto\alpha(1)\) 是环同态。其核称为增广理想。次数非负的虚特征标未必是实际特征标：在 \(C_2\) 中，若非平凡一维特征标为 \(\varepsilon\)，则 \(2\cdot1-\varepsilon\) 的次数为 \(1\)，但不可约展开有负系数。由正交关系，\(\alpha\in R(G)\) 是实际特征标，当且仅当对所有不可约 \(\chi\)，\(\langle\chi,\alpha\rangle_G\ge0\)。本章内积对第一变量共轭线性。

\subsection{特征标值的代数整数性}
\label{b5:subsec:0501:02}

\begin{definition}\label{b5:def:algebraic-integer}
	设 \(z\in\mathbb C\)。若存在 \(n\ge1\) 及整数 \(a_0,\ldots,a_{n-1}\)，使
	\[
		z^n+a_{n-1}z^{n-1}+\cdots+a_0=0,
	\]
	则称 \(z\) 为\term[daishuzhengshu]{代数整数}[algebraic integer]。
\end{definition}

下面的判据避免每遇到一个和或积便重新猜测其首一方程。
\begin{lemma}[有限生成模的整性判据]\label{b5:lem:finite-module-integrality}
	设 \(A\subset B\) 为含同一单位的交换环，\(B\) 作为 \(A\)-模由有限多个元素生成。每个 \(b\in B\) 都满足某个首一 \(A\)-系数多项式。特别地，若 \(B\subset\mathbb C\) 且 \(B\) 是有限生成的 \(\mathbb Z\)-模，则 \(B\) 中每个元素都是代数整数。
\end{lemma}
\begin{proof}
	选取模生成元 \(u_1,\ldots,u_r\)，写成 \(bu_j=\sum_i a_{ij}u_i\)，并记行向量 \(u=(u_1,\ldots,u_r)\)、矩阵 \(C=(a_{ij})\)。这些关系合在一起就是
	\[
		u(bI-C)=0.
	\]
	在交换环 \(B\) 上右乘伴随矩阵，得到
	\[
		u(bI-C)\operatorname{adj}(bI-C) =u\det(bI-C)=0.
	\]
	因此首一多项式 \(F(T)=\det(TI-C)\in A[T]\) 的值 \(F(b)\) 零化每个 \(u_j\)。由于 \(1\) 是这些生成元的 \(A\)-线性组合，\(F(b)\cdot1=0\)，故 \(F(b)=0\)。这里使用的是生成关系与伴随矩阵恒等式，不要求生成元线性无关，也没有把这组生成元当作基。
\end{proof}

\begin{proposition}\label{b5:prop:character-values-integral}
	设 \(G\) 为有限群，\(m\) 为各元素阶的最小公倍数，\(\zeta_m=e^{2\pi i/m}\)。每个复虚特征标 \(\alpha\in R(G)\) 的值均属于 \(\mathbb Z[\zeta_m]\)，并且都是代数整数。
\end{proposition}
\begin{proof}
	任一表示矩阵 \(\rho(g)\) 满足 \(\rho(g)^m=I\)。由于 \(T^m-1\) 在 \(\mathbb C\) 中无重根，\(\rho(g)\) 可对角化，其特征值都是 \(\zeta_m\) 的幂。取迹及整数线性组合便得 \(\alpha(g)\in\mathbb Z[\zeta_m]\)。环 \(\mathbb Z[\zeta_m]\) 由 \(1,\zeta_m,\ldots,\zeta_m^{m-1}\) 生成，应用\cref{b5:lem:finite-module-integrality}即可。
\end{proof}

代数整数不必是实数，更不必是通常的整数。循环群的一维特征标可以取单位根为值；整性限制的是它满足的多项式，而非其在实轴上的位置。

\subsection{表示环中的诱导法}
\label{b5:subsec:0501:03}

对有限群的子群 \(H\le G\)，限制是环同态 \(\operatorname{Res}_H^G:R(G)\rightarrow R(H)\)，诱导首先是加法群同态 \(\operatorname{Ind}_H^G:R(H)\rightarrow R(G)\)。由\cref{b5:thm:tensor-identity}，它们满足
\begin{equation}\label{b5:eq:induction-ring-projection}
	\alpha\operatorname{Ind}_H^G\beta =\operatorname{Ind}_H^G\left(\operatorname{Res}_H^G\alpha\cdot\beta\right), \qquad \alpha\in R(G),\quad\beta\in R(H).
\end{equation}
这条张量恒等式说明诱导像虽然未必含单位，却受到整个表示环乘法的控制。

\begin{proposition}\label{b5:prop:induction-ideal}
	设 \(G\) 为有限群，\(\mathcal F\) 为它的一族子群，\(R(H)\) 表示有限群 \(H\) 的复表示环。令
	\[
		I_{\mathcal F}(G)=\sum_{H\in\mathcal F}\operatorname{Ind}_H^G R(H).
	\]
	则 \(I_{\mathcal F}(G)\) 是 \(R(G)\) 的理想。因此所有虚特征标均属于这个诱导像之和，当且仅当 \(1_G\in I_{\mathcal F}(G)\)。
\end{proposition}
\begin{proof}
	定义已经给出加法子群。\eqref{b5:eq:induction-ring-projection}把一个诱导生成元与任意 \(\alpha\in R(G)\) 的乘积仍写成来自同一个 \(H\) 的诱导。若理想含 \(1_G\)，则含每个 \(\alpha=\alpha1_G\)；反向显然。
\end{proof}

诱导一般不保持乘法。若 \(H=1<G\)，则 \(\operatorname{Ind}_1^G1\) 是次数 \(|G|\) 的正则特征标。其平方的次数为 \(|G|^2\)，而 \(\operatorname{Ind}_1^G(1\cdot1)\) 的次数只有 \(|G|\)。

\subsection{分圆域与特征标的域}
\label{b5:subsec:0501:04}

\(\mathbb Q(\zeta_m)\) 称为相应的\term[fenyuanyu]{分圆域}[cyclotomic field]。对复特征标 \(\chi\)，由全部值生成的域 \(\mathbb Q\bigl(\chi(g):g\in G\bigr)\) 称为其特征标值域。\cref{b5:prop:character-values-integral}说明它包含于一个有限分圆扩张。若 \(K\subset\mathbb C\) 为子域，说一个 \(d\) 维复表示可在 \(K\) 上实现，是指存在一个 \(d\) 维 \(K\)-表示，其扩标量到 \(\mathbb C\) 后与原表示同构；等价地，原表示有一组基使全部表示矩阵的元素都属于 \(K\)。取迹说明特征标值必在 \(K\) 中，但仅知道迹在 \(K\) 中，还不能保证这样的基存在。

\ProseParagraphBreak
例如取 \(C_5=\langle c\rangle\)，令一维特征标 \(\lambda\) 满足 \(\lambda(c)=\zeta_5\)。对 \(1\le a\le4\)，有 \(\lambda(c^a)=\zeta_5^a\)，即对群元素取幂，与把特征标值中的 \(\zeta_5\) 换成 \(\zeta_5^a\) 得到同一个结果。要把这种代换用于任意特征标值，必须先证明它尊重 \(\zeta_5\) 满足的全部有理多项式关系，因而确实给出分圆域的自同构。下面的引理在任意指数 \(m\) 上建立这一点。特别地，若特征标值有理，就被这些自同构固定，对元素取与其阶互素的幂便不会改变该值；\cref{b5:prop:rational-valued-artin}将据此得到有理共轭类上的诱导表达。

\begin{lemma}\label{b5:lem:cyclotomic-power-automorphisms}
	设 \(m\ge1\) 为整数，\(\zeta_m=e^{2\pi\mathrm i/m}\)，\(a\) 为与 \(m\) 互素的整数。代换 \(\zeta_m\mapsto\zeta_m^a\) 给出 \(\mathbb Q(\zeta_m)\) 的自同构 \(\sigma_a\)。对有限群 \(G\) 的复虚特征标 \(\alpha\)，若 \(m\) 是 \(G\) 的指数，则对每个 \(g\in G\) 有
	\[
		\sigma_a\bigl(\alpha(g)\bigr)=\alpha(g^a).
	\]
\end{lemma}
\begin{proof}
	一个首一整系数多项式 \(F\) 的首一有理因子 \(f\) 仍有整数系数。确实，写 \(F=fh\)，并清除分母，把 \(f=f_0/a\)、\(h=h_0/b\) 写成本原整系数多项式除以其正的首项系数，其中 \(a,b\) 为正整数。对任一素数，\(f_0,h_0\) 的约化都是非零多项式，在有限域多项式环中其乘积仍非零，所以 \(f_0h_0\) 也是本原的。但
	\[
		f_0h_0=abF,
	\]
	而 \(F\) 首一，右侧的容度为 \(ab\)。因此 \(ab=1\)，得到 \(f=f_0\in\mathbb Z[T]\)。

	分圆多项式 \(\Phi_m(T)\) 以所有本原 \(m\) 次单位根为根；递推分解 \(T^m-1=\prod_{d\mid m}\Phi_d(T)\) 及首一多项式的整系数除法给出 \(\Phi_m\in\mathbb Z[T]\)。取 \(\zeta_m\) 的首一不可约多项式 \(f\)，由刚证的因子判据有 \(f\in\mathbb Z[T]\)，且 \(f\mid\Phi_m\)。

	若素数 \(\ell\nmid m\) 而 \(\zeta_m^\ell\) 不是 \(f\) 的根，令 \(h\ne f\) 为它在 \(\Phi_m\) 中的首一不可约因子。因为 \(h(\zeta_m^\ell)=0\)，有 \(f(T)\mid h(T^\ell)\)。这些多项式首一，整系数除法使这一整除关系可以模 \(\ell\) 约化，得到 \(\overline f\mid\overline h(T)^\ell\)。这里不要求 \(\overline f\) 不可约；它仍是正次数的首一多项式，故可选一个不可约因子 \(d\in\mathbb F_\ell[T]\)。于是 \(d\mid\overline h^\ell\)，不可约多项式在 \(\mathbb F_\ell[T]\) 中是素元，所以 \(d\mid\overline h\)。由于 \(f,h\) 是不同的不可约因子，有 \(fh\mid T^m-1\)，约化后便得
	\[
		d^2\mid\overline f\,\overline h\mid T^m-1.
	\]
	但 \(\ell\nmid m\) 时，\(T^m-1\) 与导数 \(mT^{m-1}\) 互素，不可能有这样的重因子。因此 \(\zeta_m^\ell\) 仍是 \(f\) 的根。

	若已知本原根 \(\zeta_m^b\) 是 \(f\) 的根，由 \(f\) 首一且不可约，它同样是 \(\zeta_m^b\) 在 \(\mathbb Q\) 上的最小多项式。将上述论证中的 \(\zeta_m\) 换成 \(\zeta_m^b\)，便知对每个素数 \(\ell\nmid m\)，\(\zeta_m^{b\ell}\) 仍是 \(f\) 的根。每个与 \(m\) 互素的正整数都是这些素数的乘积，\(1\) 对应空积，故全部本原 \(m\) 次单位根都是 \(f\) 的根。它们互不相同，所以
	\[
		\deg f\ge\varphi(m)=\deg\Phi_m.
	\]
	另一方面，\(f\mid\Phi_m\) 给出反向不等式；两者次数相等且均首一，因而 \(f=\Phi_m\)。

	任意与 \(m\) 互素的整数 \(a\) 都模 \(m\) 同余于一个与 \(m\) 互素的正整数，故 \(\zeta_m^a\) 也有最小多项式 \(f\)。因此代换 \(\zeta_m\mapsto\zeta_m^a\) 给出域同态，其逆由 \(a\) 模 \(m\) 的逆元给出。最后，将特征标值写成单位根的整数线性组合，逐项施行 \(\sigma_a\) 即得公式。
\end{proof}

\subsection{单位根、中心类和与标量作用的整性}
\label{b5:subsec:0501:05}

设 \(K_1,\ldots,K_r\) 为 \(G\) 的共轭类，取 \(g_j\in K_j\)。在群代数中记 \(z_j=\sum_{g\in K_j}g\)。这些\term[leihe]{类和}[class sum]构成 \(Z\bigl(\mathbb C[G]\bigr)\) 的一组基：群代数元素位于中心，恰好意味着其系数在共轭类上不变。并且
\[
	z_i z_j=\sum_k a_{ijk}z_k,\qquad a_{ijk}\in\mathbb Z_{\ge0},
\]
其中 \(a_{ijk}\) 等于把 \(K_k\) 中某个固定元素写成 \(xy\)、\(x\in K_i,y\in K_j\) 的方法数；同时共轭给出不同固定元素之间的双射。

\begin{proposition}[中心标量的整性]\label{b5:prop:central-character-integrality}
	设 \(G\) 为有限群，\(\chi\) 为次数 \(d\) 的复不可约特征标，\(g\in G\)，\(K\) 为 \(g\) 的共轭类。类和 \(z_K=\sum_{x\in K}x\in\mathbb C[G]\) 在对应不可约表示上作用为标量
	\[
		\omega_\chi(K)=\frac{|K|\cdot\chi(g)}{d},
	\]
	且此标量是代数整数。
\end{proposition}
\begin{proof}
	类和位于中心，由\cref{b5:cor:schur-algebraically-closed}，在复不可约表示上的作用为标量。取迹得到 \(d\cdot\omega_\chi(K)=|K|\cdot\chi(g)\)。整系数类和格 \(\bigoplus_j\mathbb Zz_j\) 是含单位的环，并且是有限生成的 \(\mathbb Z\)-模。\cref{b5:lem:finite-module-integrality}给出类和满足的首一整系数方程；施行表示后，标量 \(\omega_\chi(K)\) 满足同一方程。
\end{proof}

在 \(\mathbb Z[S_3]\) 中，这个过程可以直接算出。令 \(r=(123)\)，类和 \(C=r+r^{-1}\) 满足
\[
	C^2=r^2+2\cdot1+r^{-2}=C+2\cdot1.
\]
若 \(C\) 在一个复不可约表示上作用为标量 \(\omega\)，将此恒等式作用到该表示上便得 \(\omega^2-\omega-2=0\)，所以 \(\omega\in\{2,-1\}\)。整系数的类和乘法关系由此变成标量满足的首一整系数方程。对 \(S_3\) 的二维不可约特征标 \(\chi\)，\(\chi(r)=-1\)，中心标量公式给出 \(\omega=2\cdot(-1)/2=-1\)，与方程一致。

\ProseParagraphBreak
中心标量的整性比特征标值本身的整性更强，因为式中除以了次数 \(d\)。下面证明次数整除时，将先把 \(|G|/d\) 写成这些代数整数的和，再利用它是有理数而推出整数性；Burnside 证明还将把分母 \(d\) 与共轭类大小的互素条件结合。

\subsection{代数整数的基本运算与有理数判别}
\label{b5:subsec:0501:06}

\begin{lemma}\label{b5:lem:algebraic-integer-arithmetic}
	代数整数的和、差、积仍为代数整数；有理数是代数整数，当且仅当它是整数。若代数整数 \(\alpha\) 的首一最小多项式为 \(f\in\mathbb Q[T]\)，则 \(f\in\mathbb Z[T]\)，因而全部共轭根的乘积是整数。
\end{lemma}
\begin{proof}
	若 \(\alpha,\beta\) 分别满足次数 \(r,s\) 的首一整系数方程，则环 \(\mathbb Z[\alpha,\beta]\) 由 \(\alpha^i\beta^j\)、\(0\le i<r,0\le j<s\) 有限生成。\cref{b5:lem:finite-module-integrality}同时处理和、差、积。若既约分数 \(u/v\)、\(v>0\) 满足次数 \(n\) 的首一整系数方程，乘以 \(v^n\) 得 \(v\mid u^n\)，所以 \(v=1\)。

	最后，\(f\) 的每个根也是 \(\alpha\) 所满足的那个首一整系数多项式的根，因而都是代数整数。\(f\) 的系数是这些根的整数多项式，又属于 \(\mathbb Q\)，由刚证的两项结论可知均为整数。根的乘积就是带符号的常数项。
\end{proof}

\begin{corollary}[Frobenius 次数整除定理]\label{b5:cor:character-degree-divisibility}
	有限群 \(G\) 的每个复不可约特征标的次数 \(d\) 都整除 \(|G|\)。
\end{corollary}
\begin{proof}
	设该特征标为 \(\chi\)。由\cref{b5:thm:character-hom-inner-product}与\cref{b5:cor:schur-algebraically-closed}，有 \(\langle\chi,\chi\rangle_G=1\)。按共轭类求和得
	\[
		\frac{|G|}{d} =\sum_j\frac{|K_j|\cdot\chi(g_j)}{d}\,\overline{\chi(g_j)}.
	\]
	右侧各因子由\cref{b5:prop:central-character-integrality}及\cref{b5:prop:character-values-integral}均为代数整数。因此 \(|G|/d\) 是有理代数整数，由\cref{b5:lem:algebraic-integer-arithmetic}属于 \(\mathbb Z\)。
\end{proof}

整性与次数整除可参阅 Etingof 等\cite[Sections 5.2--5.3, pp. 93--98]{EtingofEtAl2011RepresentationTheory}。后文只需上述有限生成模的整性判据、类和与有理数判别，不预设一般整扩张理论。

% !TeX root = ../../Book5.tex
% 纲要标题：### 5.2 Artin 有理诱导
% 纲要位置：工作记录/计划纲要.md，第 3858 行。
% 纲要细目（写作范围）：
% * Artin 诱导定理
% * 有理系数与整数系数的区别
% * 明确所取 Artin 诱导版本：\(R_{\mathbb C}(G)\otimes_{\mathbb Z}\mathbb Q\) 由循环子群表示的诱导张成；有理系数的线性组合不意味着所有特征标值都在 \(\mathbb Q\) 中
% * 把从循环子群一般特征标诱导的表述与仅从平凡特征标诱导的有理值版本分开；用 Frobenius 互反和类函数检测组织证明

\section{Artin 有理诱导}
\label{b5:sec:0502}

每个群元素都在一个循环子群里。单靠这句话无法拼合实际表示，却足以检测类函数是否为零。借助 Frobenius 互反，把这种检测转置过来，便得到诱导的有理张成结论。

\subsection{Artin 诱导定理}
\label{b5:subsec:0502:01}

\begin{theorem}[Artin 有理诱导]\label{b5:thm:artin-induction}
	设 \(G\) 为有限群。每个复虚特征标 \(\alpha\in R(G)\) 都能写成
	\[
		\alpha=\sum_{j=1}^s a_j\operatorname{Ind}_{C_j}^G\lambda_j, \qquad a_j\in\mathbb Q,
	\]
	其中 \(C_j\le G\) 为循环子群，\(\lambda_j\) 为 \(C_j\) 的一维复特征标。等价地，这些诱导特征标张成 \(\mathbb Q\otimes_{\mathbb Z}R(G)\)。
\end{theorem}
\term[Artinyoudao]{Artin 诱导定理}[Artin's induction theorem]在这里是关于系数的结论；它没有断言 \(\alpha(g)\) 为有理数。

\begin{proof}
	先在复类函数空间中证明张成。若类函数 \(f\) 与每个 \(\operatorname{Ind}_C^G\lambda\) 正交，则由\cref{b5:cor:frobenius-character}，\(\operatorname{Res}_C^Gf\) 与循环群 \(C\) 的每个不可约特征标正交。互反等式也适用于任意类函数：\cref{b5:thm:irreducible-character-basis}把类函数写成不可约特征标的复线性组合，等式可按半线性延拓。因此 \(f|_C=0\)。对每个 \(g\in G\) 取 \(C=\langle g\rangle\)，便得 \(f(g)=0\)。有限维内积空间中正交补为零，所以这些诱导特征标复线性张成全部类函数。

	取不可约特征标为基，把全部循环子群的一维诱导特征标写为列向量，所得矩阵的元素是非负整数。刚证的结论说该矩阵满行秩，因而存在一个行数阶的非零整数子式。用对应方阵求解任一不可约特征标的坐标，逆矩阵的元素均为有理数。这就将复张成加强为有理张成，再对整数线性组合使用加法即可。
\end{proof}

最后一步必须使用整数矩阵；仅证明复线性张成，尚未自动说明一个指定向量的坐标是有理数。Etingof 等\cite[Theorem 5.26.1, Remark 5.26.2 and Corollary 5.26.3, pp. 141--142]{EtingofEtAl2011RepresentationTheory}给出了相应的子群检测表述。

\subsection{有理系数与整数系数}
\label{b5:subsec:0502:02}

\begin{corollary}\label{b5:cor:artin-finite-index}
	设 \(G\) 为有限群，\(R(G)\) 为其复表示环，令 \(I_{\mathrm{cyc}}(G)=\sum_{C\le G,\ C\text{ 循环}}\operatorname{Ind}_C^G R(C)\)。加法商群
	\[
		R(G)/I_{\mathrm{cyc}}(G)
	\]
	有限；特别地，存在同一个正整数 \(N\)，使全部 \(\alpha\in R(G)\) 均满足 \(N\alpha\in I_{\mathrm{cyc}}(G)\)。
\end{corollary}
\begin{proof}
	由\cref{b5:thm:artin-induction}，对有限多个不可约特征标分别清除有理系数的分母，再取这些分母的公倍数 \(N\)。于是 \(NR(G)\subset I_{\mathrm{cyc}}(G)\)，所述商群是有限群 \(R(G)/NR(G)\) 的商。
\end{proof}

分母确实可能无法只靠循环子群去掉。取 \(G=C_2\times C_2\)，即 Klein 四元群。\footnote{克莱因（Felix Klein，1849-1925），德国数学家，提出以变换群研究几何的埃尔朗根纲领，并研究自同构与群作用。}每个循环子群的阶为 \(1\) 或 \(2\)，从它诱导的一维特征标的次数为 \(4\) 或 \(2\)。因此循环诱导的任意整数线性组合都有偶数次数，不可能等于 \(1_G\)。若 \(H_1,H_2,H_3\) 是三个二阶子群，则
\begin{equation}\label{b5:eq:klein-artin-denominator}
	1_G=\frac12\sum_{i=1}^3\operatorname{Ind}_{H_i}^G1_{H_i} -\frac12\operatorname{Ind}_1^G1.
\end{equation}
在单位元处右侧等于 \((6-4)/2=1\)；每个非单位元恰属一个 \(H_i\)，相应置换特征标取值 \(2\)，其余项取值 \(0\)。这逐点核验了公式，也显示负系数与分母的不同职责。

\begin{proposition}\label{b5:prop:klein-cyclic-induction-index}
	设 \(G=C_2\times C_2\)，\(R(G)\) 为其复表示环，\(I_{\mathrm{cyc}}(G)\) 为从全部循环子群诱导虚特征标所成的加法子群。则
	\[
		I_{\mathrm{cyc}}(G)=\{\,\alpha\in R(G)\mid\alpha(1)\in2\mathbb Z\,\}, \qquad R(G)/I_{\mathrm{cyc}}(G)\simeq\mathbb Z/2\mathbb Z.
	\]
	特别地，\([R(G):I_{\mathrm{cyc}}(G)]=2\)。
\end{proposition}
\begin{proof}
	\(G\) 的四个不可约复特征标均一维，记为 \(e_1,e_2,e_3,e_4\)，作为 \(R(G)\) 的整数基。任意两个不同的 \(e_i,e_j\) 的商是一个非平凡一维特征标，其核 \(H\) 是二阶子群；它们在 \(H\) 上的限制相同，记为 \(\lambda\)。由 Frobenius 互反，\(\operatorname{Ind}_H^G\lambda\) 中 \(e_i,e_j\) 的重数各为一，而诱导次数为二，故
	\[
		e_i+e_j=\operatorname{Ind}_H^G\lambda\in I_{\mathrm{cyc}}(G),\qquad i\ne j.
	\]
	另一方面，每个循环子群的指数为二或四，故从它诱导的一维特征标的次数为偶数，\(I_{\mathrm{cyc}}(G)\) 中的每个元素次数也为偶数。

	为证明反向包含，对互不相同的 \(i,j,k\) 有
	\[
		e_i-e_j=(e_i+e_k)-(e_j+e_k),\qquad 2e_i=(e_i+e_j)+(e_i+e_k)-(e_j+e_k).
	\]
	所以这些差及 \(2e_i\) 都在诱导格中。若 \(\alpha=\sum_{i=1}^4n_ie_i\) 的次数 \(\sum_i n_i\) 为偶数，则
	\[
		\alpha=\sum_{i=2}^4n_i(e_i-e_1)+\Bigl(\sum_{i=1}^4n_i\Bigr)e_1
	\]
	也属于诱导格。最后，次数模二给出满同态 \(R(G)\to\mathbb Z/2\mathbb Z\)，其核恰为 \(I_{\mathrm{cyc}}(G)\)，得到所述商群和指数。
\end{proof}

\subsection{循环子群诱导与有理张成}
\label{b5:subsec:0502:03}

\begin{proposition}\label{b5:prop:subgroup-cover-induction}
	设 \(G\) 为有限群，\(\mathcal F\) 为对 \(G\)-共轭封闭的一族子群，\(R(G)\) 为复表示环，\(\operatorname{Irr}(H)\) 为 \(H\) 的不可约复特征标集合。下列条件等价：每个 \(g\in G\) 属于某个 \(H\in\mathcal F\)；所有 \(\operatorname{Ind}_H^G\theta\)、\(H\in\mathcal F\)、\(\theta\in\operatorname{Irr}(H)\) 有理张成 \(\mathbb Q\otimes_{\mathbb Z}R(G)\)。若 \(\mathcal F\) 不要求对共轭封闭，则后一条件等价于每个 \(G\)-共轭类都与某个 \(H\in\mathcal F\) 相交。
\end{proposition}
\begin{proof}
	先考虑不要求对共轭封闭的子群族。由 Frobenius 互反，与全部诱导特征标正交的类函数在每个 \(H\) 上都为零。如果每个共轭类都与这些子群相交，类函数就在全部共轭类上为零，故诱导特征标复张成全部类函数。\cref{b5:thm:artin-induction}证明中的整数矩阵秩论证给出有理张成。反之，若一个共轭类与全部 \(H\) 不相交，\cref{b5:thm:induced-character}的求和集合为空，全部诱导特征标在该类上均为零，不能张成平凡特征标。

	若子群族对共轭封闭，逐类相交就等价于逐元素覆盖：当 \(g=xhx^{-1}\)、\(h\in H\) 时，有 \(g\in xHx^{-1}\in\mathcal F\)。因此一般子群族的结论也给出所述等价条件。
\end{proof}

例如在 \(S_3\) 中只取 \(H_2=\langle(12)\rangle\)、\(H_3=\langle(123)\rangle\)。它们的并不含 \((13),(23)\)，却与每个共轭类相交。记 \(\varepsilon,\chi\) 分别为符号特征标与二维不可约特征标，\(\eta\) 为 \(H_2\) 的非平凡一维特征标，\(\lambda\) 为 \(H_3\) 的非平凡一维特征标。\Cref{b5:exm:s3-transposition-induction}给出 \(\operatorname{Ind}_{H_2}^{S_3}1=1+\chi\)、\(\operatorname{Ind}_{H_2}^{S_3}\eta=\varepsilon+\chi\)；\cref{b5:exm:s3-cyclic-induction}给出的非平凡循环诱导在三类上取值 \(2,0,-1\)，所以 \(\operatorname{Ind}_{H_3}^{S_3}\lambda=\chi\)。于是
\[
	\begin{aligned}
		1_{S_3}&=\operatorname{Ind}_{H_2}^{S_3}1_{H_2} -\operatorname{Ind}_{H_3}^{S_3}\lambda,
		\\
		\varepsilon&=\operatorname{Ind}_{H_2}^{S_3}\eta -\operatorname{Ind}_{H_3}^{S_3}\lambda.
	\end{aligned}
\]
连同 \(\chi=\operatorname{Ind}_{H_3}^{S_3}\lambda\)，这已张成 \(R(S_3)\)。这个子群族不对共轭封闭，逐类相交在这里并不等价于逐元素覆盖，这说明\cref{b5:prop:subgroup-cover-induction}中共轭封闭的条件不能省略：只有加入全部共轭子群后，逐类相交才保证逐元素覆盖。

\ProseParagraphBreak
循环子群不需要互不相交，也不需要每个元素只被覆盖一次。这里使用的是限制能检测零函数；重叠部分不会破坏正交补论证。但一般的检测只给出有理系数结果，整系数结论还需要下一节的整性论证。

\subsection{一般特征标诱导与有理值版本}
\label{b5:subsec:0502:04}

“有理”在这里涉及三个不同要求：诱导表达的系数属于 \(\mathbb Q\)，特征标在每个群元素处取值于 \(\mathbb Q\)，以及表示能在同维的 \(\mathbb Q\)-向量空间上实现。它们分别约束表达系数、函数取值和表示矩阵的定义域。\Cref{b5:thm:artin-induction}给出第一个要求；若要求仅从 \(1_C\) 诱导，得到的都是置换特征标，因而必为整数值，它们的有理线性组合也只能是有理值，涉及的是第二个要求。

\ProseParagraphBreak
例如 \(C_3=\langle c\rangle\) 的非平凡一维特征标 \(\lambda(c)=\zeta_3\) 满足 \(\lambda=\operatorname{Ind}_{C_3}^{C_3}\lambda\)，诱导表达的系数为 \(1\)，但它的值不是有理数，因而不能由平凡特征标的循环诱导有理张成。若一个复表示可在 \(\mathbb Q\) 上实现，其特征标当然为有理值；反向是否成立，则须进一步检查同维的有理矩阵能否满足群的关系。

\begin{proposition}\label{b5:prop:rational-values-no-real-model}
	设 \(Q_8=\{\pm1,\pm\mathbf i,\pm\mathbf j,\pm\mathbf k\}\) 为四元数群，\(\theta\) 为它的二维复不可约特征标，其值为
	\[
		\theta(1)=2,\qquad\theta(-1)=-2,\qquad \theta(g)=0\quad(g\notin\{1,-1\}).
	\]
	虽然 \(\theta\) 是整数值特征标，但不存在特征标为 \(\theta\) 的二维实表示，因而相应复表示也不能在 \(\mathbb Q\) 上实现。
\end{proposition}
\begin{proof}
	假设存在这样的二维实表示。群元素 \(-1\) 的表示矩阵平方为 \(I\)，在实数域上可对角化，特征值为 \(1\) 或 \(-1\)；迹为 \(\theta(-1)=-2\) 迫使它等于 \(-I\)。因此 \(\mathbf i,\mathbf j\) 的实矩阵 \(A,B\) 满足
	\[
		A^2=B^2=-I,\qquad AB=-BA.
	\]
	取非零实向量 \(v\)。若 \(Av=tv\)、\(t\in\mathbb R\)，则 \(-v=A^2v=t^2v\)，不可能，故 \(v,Av\) 构成一组实基。在这组基下，反交换关系给出
	\[
		A=\begin{pmatrix}0&-1\\1&0\end{pmatrix},\qquad
		B=\begin{pmatrix}a&b\\b&-a\end{pmatrix},\qquad a,b\in\mathbb R.
	\]
	但 \(B^2=(a^2+b^2)I\)，不可能等于 \(-I\)。二维有理表示扩标量到 \(\mathbb R\) 会给出同维实表示，所以有理模型同样不存在。
\end{proof}

这个反例的特征标值域为 \(\mathbb Q\)，表示矩阵仍无法在该域中实现。\Cref{b5:thm:artin-induction}仍给出 \(\theta\) 的有理系数诱导表达；所以表达系数有理、函数取值有理，并不能代替同维表示的有理实现。

\ProseParagraphBreak
有理值虚特征标受到的幂运算约束由\cref{b5:lem:cyclotomic-power-automorphisms}给出。对有限群 \(G\) 中的元素 \(u,v\)，若存在与 \(|u|\) 互素的整数 \(a\)，使 \(v\) 共轭于 \(u^a\)，则称 \(u,v\) 有理共轭；这等价于它们生成的循环子群共轭。

\begin{proposition}[有理值版本]\label{b5:prop:rational-valued-artin}
	设 \(G\) 为有限群。所有在 \(G\) 的有理共轭类上常值的函数 \(f:G\to\mathbb Q\)，由 \(\operatorname{Ind}_C^G1_C\) 在 \(\mathbb Q\) 上张成，其中 \(C\le G\) 遍历循环子群，\(1_C\) 为平凡复特征标。特别地，每个取值于 \(\mathbb Q\) 的复虚特征标都有这种表达。
\end{proposition}
\begin{proof}
	从循环子群共轭类中选代表 \(C_1,\ldots,C_s\)，各取生成元 \(c_i\)，按 \(|C_i|\) 从大到小排列。置换特征标 \(\pi_j=\operatorname{Ind}_{C_j}^G1\) 在 \(c_i\) 处非零，恰好在某个共轭的 \(C_j\) 包含 \(C_i\) 时发生。因此矩阵 \(\bigl(\pi_j(c_i)\bigr)\) 为三角矩阵，同阶的不同代表之间该值为零。对角元计数的是被 \(c_i\) 固定的左陪集 \(xC_i\)，其条件为 \(x^{-1}c_ix\in C_i\)。由于 \(c_i\) 生成整个 \(C_i\)，这个条件等价于 \(x^{-1}C_ix=C_i\)，即 \(x\in N_G(C_i)\)。每个陪集有 \(|C_i|\) 个代表，所以
	\[
		\pi_i(c_i)=\frac{|N_G(C_i)|}{|C_i|}>0.
	\]
	这是正规化子的指数，而不是中心化子的指数：例如 \(C_3\trianglelefteq S_3\) 时，正规化子为 \(S_3\)，计数为 \(2\)；中心化子只有 \(C_3\)，若误用它就会得到 \(1\)。矩阵为整系数且可逆，故每一组有理数值都能由这些 \(\pi_j\) 的有理线性组合得到。一个有理共轭类恰对应一个循环子群共轭类，所以这已覆盖所述函数空间。

	若虚特征标 \(\alpha\) 为有理值，由\cref{b5:lem:cyclotomic-power-automorphisms}得 \(\alpha(g^a)=\alpha(g)\)，其中 \(a\) 与群指数互素。若只给定 \(a\) 与 \(|g|\) 互素，可在同一模 \(|g|\) 剩余类中选择与群指数互素的整数：对未整除 \(|g|\) 的素因子另规定非零剩余类，再用中国剩余定理。于是 \(\alpha\) 在有理共轭类上常值，应用前一结论即可。
\end{proof}

循环子群的包含关系可以在 \(G=C_6=\langle c\rangle\) 中直接看见。按阶从大到小列出 \(C_6,C_3=\langle c^2\rangle,C_2=\langle c^3\rangle,C_1=1\)，生成元依次取 \(c,c^2,c^3,1\)。因为 \(G\) 交换，\((\operatorname{Ind}_H^G1_H)(g)\) 在 \(g\in H\) 时等于 \([G:H]\)，否则为零；所以上述求值矩阵为
\[
	M=\begin{pmatrix}
		1&0&0&0\\
		1&2&0&0\\
		1&0&3&0\\
		1&2&3&6
	\end{pmatrix},\qquad \det M=36.
\]
每一行的生成元只能属于包含它所生成的循环子群的那些子群，因而严格上三角部分全为零。非零对角元使系数可以从第一行起逐行解出；求解时须分别除以 \(1,2,3,6\)，因此所得系数属于 \(\mathbb Q\)，未必为整数。例如在四个有理共轭类上依次取值 \(0,1,0,0\) 的函数满足
\[
	f=\frac12\operatorname{Ind}_{C_3}^{C_6}1_{C_3} -\frac16\operatorname{Ind}_{1}^{C_6}1,
\]
矩阵可逆保证这个表达在所列四个诱导特征标中唯一。这里讨论的 \(f\) 是一般有理值函数；命题并未要求它是虚特征标。

% !TeX root = ../../Book5.tex
% 纲要标题：### 5.3 Brauer 诱导
% 纲要位置：工作记录/计划纲要.md，第 3865 行。
% 纲要细目（写作范围）：
% * 初等子群及一维特征标
% * Brauer 诱导定理的明确版本
% * 与原纲要未指明版本的“Brauer 定理”对接
% * 定义 \(p\)-初等群为 \(P\times C\)，其中 \(P\) 为有限 \(p\)-群、\(C\) 为阶与 \(p\) 互素的循环群；所有素数一起使用
% * 目标版本为复虚特征标是从初等子群诱导的特征标的整数线性组合，并细化到适当子群的一维特征标；完整整性证明在第5.4节
% * 本章不以前置第8章的模表示、Brauer 特征标或分解数为条件

\section{Brauer 诱导}
\label{b5:sec:0503}

Klein 四元群的例子说明，循环子群有时太小，不能整系数地生成表示环。需要加入的不是任意复杂的子群，而是把一个素数幂阶部分与一个互素的循环部分并列起来。本节确定目标及其准确范围，证明集中在下一节。

\subsection{初等子群与一维特征标}
\label{b5:subsec:0503:01}

\begin{definition}\label{b5:def:elementary-group}
	设 \(p\) 为素数，\(E\) 为有限群。若 \(E\) 同构于 \(P\times C\)，其中 \(P\) 是有限 \(p\)-群，\(C\) 是阶与 \(p\) 互素的循环群，则称 \(E\) 为\term[pchudengqun]{\(p\)-初等群}[\(p\)-elementary group]。若存在某个素数 \(p\) 使此条件成立，则称 \(E\) 为\term[chudengqun]{初等群}[elementary group]。有限群 \(G\) 中具有上述结构的子群称为相应的初等子群。
\end{definition}

允许 \(P=1\) 或 \(C=1\)。因此全部有限 \(p\)-群都是 \(p\)-初等群；每个循环群对每一个素数 \(p\) 都是 \(p\)-初等群，只需把其阶分成 \(p\) 部分与互素部分。初等群不一定交换，例如 \(Q_8\times C_3\) 是 \(2\)-初等群。

\ProseParagraphBreak
Klein 四元群 \(G=C_2\times C_2\) 的循环诱导子格指数恰为 \(2\)，见\cref{b5:prop:klein-cyclic-induction-index}。若允许 \(2\)-初等子群，就能取 \(P=G\)、\(C=1\)，从全群自身的恒等诱导得到全部特征标。原来循环子群不能处理的 \(2\) 部分因此得以消去。对奇素数 \(p\)，\(G\) 的 \(p\)-初等子群仍只有循环子群，指数仍为 \(2\)，但在反演与 \(p\) 互素整数得到的 \(\mathbb Z_{(p)}\) 中，这个指数已可逆。\cref{b5:prop:klein-elementary-local-indices}将给出具体比较。这说明固定素数 \(p\) 时，要排除的是指数中的 \(p\) 部分；允许一个 \(p\)-群因子，正是为了处理循环子群不足以消去的这一部分。

\ProseParagraphBreak
从一维特征标 \(\lambda:H\rightarrow\mathbb C^\times\) 诱导得到的特征标称为\term[danxiangshitezhengbiao]{单项式特征标}[monomial character]。取左陪集代表后，诱导表示的每个群元素把每个基向量送到另一基向量的非零标量倍，所以矩阵每行、每列恰有一个非零元。名称来自这个具体矩阵形状，并不意味着诱导结果总不可约。

\subsection{Brauer 诱导定理的整系数版本}
\label{b5:subsec:0503:02}

\begin{theorem}[Brauer 整诱导]\label{b5:thm:brauer-induction}
	设 \(G\) 为有限群。每个复虚特征标 \(\alpha\in R(G)\) 都是从初等子群的复特征标诱导而得的特征标的整数线性组合。进一步，存在初等子群 \(H_1,\ldots,H_t\le G\)、它们的一维复特征标 \(\lambda_i\) 及整数 \(n_i\)，使
	\[
		\alpha=\sum_{i=1}^t n_i\operatorname{Ind}_{H_i}^G\lambda_i.
	\]
\end{theorem}
这就是本章所用的\term[Braueryoudao]{Brauer 诱导定理}[Brauer's induction theorem]。整系数结论将在\cref{b5:prop:brauer-integral-image}证明，再由\cref{b5:prop:elementary-monomial}化为一维特征标的诱导。

\ProseParagraphBreak
这是表示环中的加法恒等式。若 \(\alpha\) 是实际表示的特征标，右侧仍可能有负系数；把正负部分移到两边才得到实际表示的同构。定理也未声称每个不可约表示单独就是一个一维表示的诱导。

\subsection{Brauer 诱导定理的版本辨析}
\label{b5:subsec:0503:03}

Artin 定理允许循环子群、一般一维特征标和有理系数；Brauer 定理扩大子群族，以此把系数收紧为整数。如果只保留平凡特征标，则所有诱导都是整数值函数，不能产生 \(C_3\) 的非实一维特征标。故“只从平凡表示诱导”不能作为这里的 Brauer 定理。

\ProseParagraphBreak
初等群 \(P\times C\) 与初等交换 \(p\)-群 \(C_p^r\) 是不同概念：前者允许非交换 \(P\) 和互素循环因子。把直积放宽为某些半直积，会得到别的诱导定理中使用的子群族；本章始终使用上述直积定义。Serre 的书中将相应整诱导集中放在 Chapter 10，可作为进一步阅读\cite[Chapter 10, pp. 74--80]{Serre1977FiniteRepresentations}。

\subsection{\texorpdfstring{\(p\)-初等群的内部判据与子群结构}{p-初等群的内部判据与子群结构}}
\label{b5:subsec:0503:04}

定义中的直积具有可操作的内部判据：\(E\) 中存在子群 \(P,C\)，满足 \(E=PC\)、\(P\cap C=1\)、每个 \(P\) 元素与每个 \(C\) 元素交换，且阶与循环性满足\cref{b5:def:elementary-group}。此时 \((x,c)\mapsto xc\) 是群同构。

\begin{proposition}\label{b5:prop:elementary-subgroups}
	设 \(p\) 为素数，\(P\) 为有限 \(p\)-群，\(C\) 为阶与 \(p\) 互素的有限循环群，\(E=P\times C\)。每个子群 \(H\le E\) 都是 \(p\)-初等群，具体地
	\[
		H=(H\cap P)\times(H\cap C),
	\]
	其中把 \(P,C\) 视作 \(E\) 的两个自然子群。
\end{proposition}
\begin{proof}
	对 \(h=xc\in H\)，\(x\in P,c\in C\)，元素 \(x,c\) 的阶互素。中国剩余定理给出整数 \(a\)，使 \(a\equiv1\pmod{|x|}\)、\(a\equiv0\pmod{|c|}\)，所以 \(h^a=x\in H\)，继而 \(c=x^{-1}h\in H\)。因此 \(H=(H\cap P)(H\cap C)\)。两因子交换且交为 \(1\)，第二个因子是循环群的子群，阶仍与 \(p\) 互素。
\end{proof}

例如 \(S_3\) 不是初等群：其中心为 \(1\)，且它的阶不是素数幂；若它是 \(P\times C\) 且 \(C\ne1\)，则 \(C\) 必在中心中，矛盾。不过其二阶、三阶子群都是初等子群，已经足以给出整诱导。

\subsection{从初等子群到一维特征标的诱导}
\label{b5:subsec:0503:05}

设 \(E=P\times C\)。对复不可约 \(E\)-表示 \(V\)，由\cref{b5:cor:schur-algebraically-closed}，中心中的 \(C\) 作用为一个一维特征标 \(\lambda\)。任何 \(P\)-不变子空间自动也被这个标量作用保持，所以 \(V|_P\) 不可约。故 \(V\) 的特征标为 \((x,c)\mapsto\theta(x)\lambda(c)\)，其中 \(\theta\in\operatorname{Irr}(P)\)。

\ProseParagraphBreak
剩下的任务是证明 \(p\)-群的不可约特征标可从较小子群的一维特征标诱导。一旦写出 \(\theta=\operatorname{Ind}_Q^P\mu\)，就有
\[
	\theta\otimes\lambda =\operatorname{Ind}_{Q\times C}^{P\times C}(\mu\otimes\lambda).
\]
右侧子群仍是 \(p\)-初等群。这一步所需的 \(p\)-群结论将在\cref{b5:thm:p-group-monomial}证明；再由诱导传递性与从 \(E\) 到 \(G\) 的诱导复合。

\subsection{Brauer 诱导与模表示的关系}
\label{b5:subsec:0503:06}

本章始终讨论复表示。下一节会把特征标值模某个素理想约化到特征 \(p\) 的域中，但这只是对一个交换数环取商，并不把表示空间或群代数改为特征 \(p\) 上的对象。模表示在 \(p\,\bigm\vert\,|G|\) 时通常不完全可约。

\ProseParagraphBreak
第八章的 Brauer 特征标只在 \(p\)-正则元素上定义，并处理特征 \(p\) 的表示；这里的 Brauer 诱导则是全部群元素上的普通复特征标恒等式。名字相同，所研究的对象却不同。

% !TeX root = ../../Book5.tex
% 纲要标题：### 5.4 初等子群检测与 Brauer 整诱导的证明
% 纲要位置：工作记录/计划纲要.md，第 3874 行。
% 纲要细目（写作范围）：
% * 把诱导像组织为复表示环中的加法子群，利用限制、诱导与张量恒等式控制其环论行为
% * 建立初等子群上的限制检测与整性引理，分素数处理局部化后的整数系数问题，再由有限生成交换群的局部检测回到整数系数
% * 分别完成“存在某个整数倍的诱导表达”与“无需整体分母的整诱导表达”，指出第二步是超出 Artin 有理诱导的实质
% * 解释初等群上的表示如何进一步化为从适当子群一维表示的诱导；用完整权空间建立惯性群表示与诱导，不把权空间误当成一维；逐项说明所用子群仍处于允许的初等范围
% * 用小群的虚特征标表达说明整数系数可以为负，不把虚表示恒等式误解为实际表示的直接和分解

\section{初等子群检测与 Brauer 整诱导的证明}
\label{b5:sec:0504}

Artin 定理已经给出有限指数的子格。现在必须证明这个指数没有任何素因子。对每个素数 \(p\)，我们反演所有与 \(p\) 互素的整数，证明 \(p\)-初等子群的诱导在这个局部化以后生成整个表示环，从而排除指数中的素因子 \(p\)。论证中所需的交换环知识将在使用处给出证明。

\ProseParagraphBreak
张量恒等式说明诱导像是理想。为了证明它在 \(p\)-局部化后等于全环，我们把局部化表示环 \(R(G)_{(p)}\) 的极大理想写成模 \(p\) 的特征标值检测，并将检测元素约化为 \(p\)-正则部分。在相应中心化子中构造 \(p\)-初等子群，就能找到不落入该极大理想的诱导特征标。对每个素数重复这一论证，才得到整系数等式。

\subsection{诱导像与复表示环}
\label{b5:subsec:0504:01}

固定有限群 \(G\) 与素数 \(p\)，记 \(\mathcal E_p(G)\) 为全部 \(p\)-初等子群，并令
\[
	I_p(G)=\sum_{E\in\mathcal E_p(G)}\operatorname{Ind}_E^G R(E).
\]
由\cref{b5:prop:induction-ideal}，这是 \(R(G)\) 的理想。每个循环子群都是 \(p\)-初等群，因此\cref{b5:cor:artin-finite-index}给出
\begin{equation}\label{b5:eq:artin-elementary-finite-index}
	I_{\mathrm{cyc}}(G)\subset I_p(G)\subset R(G), \qquad \bigl[R(G):I_p(G)\bigr]<\infty.
\end{equation}
还需排除这个有限指数中的素因子 \(p\)。

\ProseParagraphBreak
记 \(\mathbb Z_{(p)}=\{a/b\in\mathbb Q:p\nmid b\}\)。对加法群 \(M\)，记 \(M_{(p)}=\mathbb Z_{(p)}\otimes_{\mathbb Z}M\)，可将元素写成 \(x/s\)，其中 \(p\nmid s\)；判据 \(x/1=0\) 等价于存在这样的 \(s\) 使 \(sx=0\)。我们将证明
\[
	I_p(G)_{(p)}=R(G)_{(p)}.
\]
这并不要求单独固定一个 \(p\) 就能在整数上生成 \(R(G)\)；它只说明最终的障碍不含 \(p\) 部分。

\begin{proposition}\label{b5:prop:klein-elementary-local-indices}
	设 \(G=C_2\times C_2\) 为 Klein 四元群，\(p\) 为素数，\(I_p(G)\) 为其全部 \(p\)-初等子群的复表示环诱导像之和，\(I_{\mathrm{cyc}}(G)\) 为循环子群的诱导像之和。则
	\[
		\begin{aligned}
			p=2&:\quad I_2(G)=R(G),\quad [R(G):I_2(G)]=1;\\
			p\ne2&:\quad I_p(G)=I_{\mathrm{cyc}}(G),\quad [R(G):I_p(G)]=2.
		\end{aligned}
	\]
	对每个素数 \(p\)，均有 \(I_p(G)_{(p)}=R(G)_{(p)}\)，其中下标 \((p)\) 表示反演所有与 \(p\) 互素的整数。
\end{proposition}
\begin{proof}
	当 \(p=2\) 时，\(G\) 自身是 \(2\)-初等群，故 \(\operatorname{Ind}_G^G R(G)=R(G)\) 已包含在 \(I_2(G)\) 中。当 \(p\ne2\) 时，\(G\) 的任何子群都没有非平凡 \(p\)-子群，所以 \(p\)-初等子群中的 \(P\) 因子只能为 \(1\)，这些子群恰为循环子群。\cref{b5:prop:klein-cyclic-induction-index}遂给出指数 \(2\)。其证明中的偶数总系数格描述还表明 \(2R(G)\subset I_p(G)\)；因 \(2\) 在 \(\mathbb Z_{(p)}\) 中可逆，局部化后任一 \(\alpha\) 都可写成 \((2\alpha)/2\)，故局部化诱导像等于全环。\(p=2\) 时则无需反演分母。
\end{proof}

对这个群，扩大到 \(2\)-初等子群消去了指数中的 \(2\) 因子；固定奇素数时并未在整数上消去它，而是借局部化将它变为单位。下面的一般证明要对每个群完成相同的逐素数排除。

\subsection{限制检测、局部化与整性引理}
\label{b5:subsec:0504:02}

需要先把 \(p\)-局部化表示环 \(R(G)_{(p)}\) 的极大理想解释为模 \(p\) 的特征标值检测。为此，先证明有限交换环扩张中的极大理想提升引理。

\begin{lemma}[有限扩张中的极大理想提升]\label{b5:lem:finite-extension-maximal-lift}
	设 \(A\subset B\) 为含同一单位的交换环，且 \(B\) 为有限生成 \(A\)-模。对每个极大理想 \(\mathfrak m\subset A\)，存在极大理想 \(\mathfrak n\subset B\) 使 \(\mathfrak n\cap A=\mathfrak m\)。
\end{lemma}
\begin{proof}
	令 \(S=A\setminus\mathfrak m\)。分式环 \(A_{\mathfrak m}=S^{-1}A\) 的非单位恰组成理想 \(\mathfrak mA_{\mathfrak m}\)，且 \(A_{\mathfrak m}\subset S^{-1}B\) 仍为单射。因此 \(S^{-1}B\ne0\)，并为有限生成 \(A_{\mathfrak m}\)-模。

	若 \(\mathfrak mB=B\)，则有限生成模 \(M=S^{-1}B\) 满足 \(\mathfrak mA_{\mathfrak m}M=M\)。取模生成元 \(u_1,\ldots,u_r\)，写成 \(u_j=\sum_i c_{ij}u_i\)，其中全部 \(c_{ij}\in\mathfrak mA_{\mathfrak m}\)。记 \(u=(u_1,\ldots,u_r)\)、\(C=(c_{ij})\)，这些关系给出 \(u(I-C)=0\)。右乘伴随矩阵便得
	\[
		u(I-C)\operatorname{adj}(I-C)=u\det(I-C)=0.
	\]
	因 \(u_j\) 生成 \(M\)，这个行列式零化整个 \(M\)，无须假定生成元线性无关。它模 \(\mathfrak mA_{\mathfrak m}\) 为 \(1\)，故为单位，迫使 \(M=0\)，矛盾。因此 \(\mathfrak mB\ne B\)。取包含它的极大理想 \(\mathfrak n\)。其收缩是包含 \(\mathfrak m\) 的真理想，只能等于 \(\mathfrak m\)。
\end{proof}

存在极大理想只需通常的链并与 Zorn 引理\footnote{佐恩（Max August Zorn，1906-1993），德国数学家，提出与选择公理等价的极大原理，并研究代数结构。}：真理想的链之并仍不含 \(1\)。在本节中 \(B/\mathfrak mB\) 还是域 \(A/\mathfrak m\) 上的有限维非零代数，也可直接从真理想的维数中取最大者。

\ProseParagraphBreak
求值映射的目标环须同时容纳特征标值并保留模 \(p\) 的信息。前一个要求不能只靠 \(\mathbb Z_{(p)}\) 满足：例如 \(C_3\) 的非平凡一维复特征标在生成元处取值 \(\zeta_3\)，并不属于这个有理数环。\cref{b5:prop:character-values-integral}则说明，若 \(m\) 是 \(G\) 的指数，全部值都落在 \(O=\mathbb Z[\zeta_m]\) 中。若把目标扩大为 \(\mathbb C\)，素数 \(p\) 已可逆，\(p\mathbb C=\mathbb C\)，因而 \(\mathbb Z_{(p)}\) 的极大理想 \((p)\) 无法从 \(\mathbb C\) 的极大理想收缩得到。\(O_{(p)}\) 仍是有限生成 \(\mathbb Z_{(p)}\)-模，正好允许使用刚证的提升引理。取它在各共轭类上的直积，便能把表示环的极大理想变成某个共轭类上的模 \(p\) 求值条件。

\begin{lemma}[极大理想的求值描述]\label{b5:lem:representation-maximal-evaluation}
	设 \(G\) 为有限群，\(m\) 为其指数，\(\zeta_m=e^{2\pi\mathrm i/m}\)，\(O=\mathbb Z[\zeta_m]\)，\(p\) 为素数，\(R(G)\) 为复表示环。下标 \((p)\) 表示反演所有与 \(p\) 互素整数的局部化。对 \(R(G)_{(p)}\) 的每个极大理想 \(\mathfrak m\)，存在 \(g\in G\) 及 \(O_{(p)}\) 的极大理想 \(\mathfrak q\)，使剩余域 \(O_{(p)}/\mathfrak q\) 的特征为 \(p\)，且
	\[
		\mathfrak m=\bigl\{\,\alpha\in R(G)_{(p)}:\alpha(g)\in\mathfrak q\,\bigr\}.
	\]
\end{lemma}
\begin{proof}
	设共轭类数为 \(r\)。在每个共轭类代表处取值给出单射环同态
	\[
		R(G)_{(p)}\longrightarrow B=O_{(p)}^r.
	\]
	单射性来自特征标作为类函数的相等判据；由\cref{b5:prop:character-values-integral}，各分量的特征标值都属于 \(O_{(p)}\)，故这确实是到 \(B\) 的求值映射。\(O\) 为有限生成 \(\mathbb Z\)-模，所以 \(B\) 为有限生成 \(\mathbb Z_{(p)}\)-模。在这个嵌入下，\(\mathbb Z_{(p)}\) 通过 \(a\mapsto a\,1_G\) 嵌入 \(R(G)_{(p)}\)，其在 \(B\) 中的像是对角标量 \((a,\ldots,a)\)。因此一组 \(\mathbb Z_{(p)}\)-模生成元也生成 \(B\) 作为 \(R(G)_{(p)}\)-模。\cref{b5:lem:finite-extension-maximal-lift}遂把 \(\mathfrak m\) 提升为 \(B\) 的极大理想。

	有限直积的极大理想只在一个分量上取极大理想：各分量的单位幂等元在域商中只能取 \(0\) 或 \(1\)，它们的和是 \(1\)，且彼此乘积为零，因此恰有一个取 \(1\)。这便给出某个代表 \(g\) 及 \(O_{(p)}\) 的极大理想 \(\mathfrak q\)。

	最后，\(O\) 是有限生成 \(\mathbb Z\)-模，局部化再取商后，域 \(O_{(p)}/\mathfrak q\) 仍是有限生成 \(\mathbb Z_{(p)}\)-模。由\cref{b5:lem:finite-module-integrality}，它整于 \(\mathbb Z_{(p)}\) 在其中的像。若域 \(L\) 整于子整环 \(D\)，则 \(D\) 本身是域：对 \(0\ne a\in D\)，把 \(a^{-1}\in L\) 的首一整方程乘以适当的 \(a\) 的幂，就得到 \(a^{-1}\in D\)。因此 \(\mathfrak q\cap\mathbb Z_{(p)}\) 为极大理想，而 \(\mathbb Z_{(p)}\) 中不属于 \(p\mathbb Z_{(p)}\) 的元素皆可逆，故这是它唯一的极大理想，只能有 \(\mathfrak q\cap\mathbb Z_{(p)}=p\mathbb Z_{(p)}\)。剩余域的特征遂为 \(p\)。
\end{proof}

群元素 \(g\) 有唯一分解 \(g=su=us\)，其中 \(s,u\) 均为 \(g\) 的幂，\(|s|\) 与 \(p\) 互素，\(|u|\) 为 \(p\) 的幂。这由循环群 \(\langle g\rangle\) 的阶的互素分解得到。这样的 \(s\) 称为 \(g\) 的 \(p\)-正则部分；阶与 \(p\) 互素的元素称为\term[pzhengzeyuansu]{\(p\)-正则元素}[\(p\)-regular element]。

\ProseParagraphBreak
表示环的求值像在整数上通常不是任意坐标的直积。以 \(C_2\) 为例，若非平凡一维特征标记为 \(t\)，则 \(t^2=1\)，而 \(1,t\) 是表示环的整数基，所以
\[
	R(C_2)=\mathbb Z[t]/(t^2-1).
\]
在单位元与非单位元处取值，把 \(a+bt\) 送到 \((a+b,a-b)\)。像恰为同奇偶的整数对，因为反解是 \(a=(u+v)/2\)、\(b=(u-v)/2\)。局部化到奇素数 \(p\) 时，\(2\) 可逆，这个求值映射成为
\[
	R(C_2)_{(p)}\cong\mathbb Z_{(p)}\times\mathbb Z_{(p)},
\]
两个极大理想分别由 \(t=1\) 与 \(t=-1\) 的模 \(p\) 求值给出。局部化到 \(p=2\) 时，模 \(2\) 后却有 \(t^2-1=(t-1)^2\)，两个求值都变成 \(t\mapsto1\)，得到同一个极大理想 \((2,t-1)\)。由刚证的求值描述，它也是此时唯一的极大理想。于是原本不同的两个共轭类在模 \(2\) 的检测中合并，不能任意指定它们的两个剩余值。

\begin{lemma}\label{b5:lem:character-p-regular-reduction}
	设 \(G\) 为有限群，\(p\) 为素数，\(m\) 为 \(G\) 的指数，\(\zeta_m=e^{2\pi\mathrm i/m}\)，\(O=\mathbb Z[\zeta_m]\)，\(R(G)\) 为复表示环，下标 \((p)\) 表示反演所有与 \(p\) 互素整数的局部化。设 \(\mathfrak q\subset O_{(p)}\) 为剩余特征是 \(p\) 的极大理想。若 \(g\in G\) 有分解 \(g=su=us\)，其中 \(s,u\) 均为 \(g\) 的幂、\(|s|\) 与 \(p\) 互素、\(|u|\) 为 \(p\) 的幂，则每个 \(\alpha\in R(G)_{(p)}\) 均满足
	\[
		\alpha(g)\equiv\alpha(s)\pmod{\mathfrak q}.
	\]
\end{lemma}
\begin{proof}
	对实际表示，交换且可对角化的矩阵 \(\rho(s),\rho(u)\) 可同时对角化：先按 \(\rho(s)\) 的特征空间分解，再在每个被 \(\rho(u)\) 保持的空间内对角化。共同特征向量上的特征值分别记作 \(\xi,\eta\)，其中 \(\eta^{p^a}=1\)。在特征 \(p\) 的剩余域中，\(T^{p^a}-1=(T-1)^{p^a}\)，故 \(\eta\equiv1\)。于是 \(\xi\eta\equiv\xi\)，求迹即得结论。再取整数线性组合并除以与 \(p\) 互素的分母，便覆盖全部 \(\alpha\)。
\end{proof}

取 \(G=C_6=\langle c\rangle\)、\(p=2\)，对 \(g=c\) 有 \(s=c^4\)、\(u=c^3\)，两者的阶分别为 \(3,2\)，且 \(su=c\)。令一维复特征标 \(\lambda\) 满足 \(\lambda(c)=\zeta_6=e^{2\pi\mathrm i/6}\)，取 \(O=\mathbb Z[\zeta_6]\)，并取 \(O_{(2)}\) 的剩余特征为 \(2\) 的极大理想 \(\mathfrak q\)。此时
\[
	\lambda(s)=\zeta_6^4,\qquad \lambda(u)=-1,\qquad \lambda(g)=-\lambda(s).
\]
在 \(O_{(2)}\) 的剩余特征为 \(2\) 的域商中，\(-1\) 变成 \(1\)，所以 \(\lambda(g)\equiv\lambda(s)\pmod{\mathfrak q}\)。三阶因子却仍非平凡：\(\zeta_6^4\) 满足 \(T^2+T+1=0\)，而此多项式在特征 \(2\) 时不以 \(1\) 为根，其像仍有阶 \(3\)。这里消去的是数环商中二阶单位根的因子；群元素的分解与原来的复表示空间仍是上述对象。

\subsection{从有理诱导到整系数诱导}
\label{b5:subsec:0504:03}

\begin{theorem}[逐素数的整诱导]\label{b5:thm:brauer-local-induction}
	设 \(G\) 为有限群，\(p\) 为素数，\(R(G)\) 为复表示环，令 \(I_p(G)=\sum_{E\le G,\ E\text{ 为 }p\text{-初等}}\operatorname{Ind}_E^G R(E)\)。下标 \((p)\) 表示反演所有与 \(p\) 互素整数的局部化，则
	\[
		I_p(G)_{(p)}=R(G)_{(p)}.
	\]
	因此有限指数 \(\bigl[R(G):I_p(G)\bigr]\) 与 \(p\) 互素。
\end{theorem}
\begin{proof}
	若理想 \(I_p(G)_{(p)}\) 为真理想，则它包含在某个极大理想 \(\mathfrak m\) 中。由\cref{b5:lem:representation-maximal-evaluation}及\cref{b5:lem:character-p-regular-reduction}，存在 \(p\)-正则元素 \(s\) 及剩余特征为 \(p\) 的 \(\mathfrak q\subset O_{(p)}\)，使
	\[
		\alpha\in\mathfrak m\quad\Longleftrightarrow\quad \alpha(s)\equiv0\pmod{\mathfrak q}.
	\]
	我们构造一个 \(p\)-初等诱导，使其在 \(s\) 处不满足此同余。

	令 \(C=\langle s\rangle\)，在中心化子 \(C_G(s)\) 中取 Sylow \(p\)-子群 \(P\)。由于 \(C\) 在该中心化子内居中，且 \(|C|\) 与 \(|P|\) 互素，
	\[
		E=CP=C\times P
	\]
	为 \(p\)-初等子群。写 \(|C_G(s)|=p^a b\)、\(p\nmid b\)，则 \(|P|=p^a\)，且 \(|C|\,\bigm\vert\,b\)，所以
	\[
		d=\bigl[C_G(s):E\bigr]=\frac{b}{|C|}
	\]
	与 \(p\) 互素。中心化子的选择同时保证两件事：\(P\) 与 \(C\) 逐元素交换，且 \(P\) 含有中心化子的全部 \(p\) 部分。

	正规化子只保证共轭保持 \(C\)，不足以保证这两个因子交换。例如 \(G=S_3\)、\(p=2\)、\(s=(123)\) 时，\(N_G(C)=S_3\)，其中换位把 \(s\) 共轭为 \(s^{-1}\)。取正规化子中的二阶子群会得到非平凡半直积 \(S_3\)，而不是 \(C_3\times C_2\)。中心化子却是 \(C_G(s)=C_3\)，它的 Sylow \(2\)-子群为 \(1\)，因而上述构造给出 \(E=C_3\)、\(d=1\)。

	对 \(C\) 的一维特征标 \(\lambda\)，将其延拓为 \(E\) 上在 \(P\) 处平凡的特征标 \(\widetilde\lambda\)。诱导公式\cref{b5:thm:induced-character}给出
	\begin{equation}\label{b5:eq:elementary-induction-evaluation}
		\left(\operatorname{Ind}_E^G\widetilde\lambda\right)(s) =d\sum_{t\in s^G\cap C}\lambda(t).
	\end{equation}
	详细地说，\(E\) 中与 \(s\) 共轭的元素阶与 \(p\) 互素，故其 \(P\) 分量为 \(1\)，恰落在 \(C\) 内。对每个固定的 \(t\in s^G\cap C\)，满足 \(x^{-1}sx=t\) 的 \(x\in G\) 有 \(|C_G(s)|\) 个，除以 \(|E|\) 即得上式。

	集合 \(T=s^G\cap C\) 非空。平凡特征标给出的和是 \(|T|\)，若 \(p\,\bigm\vert\,|T|\)，这个值模 \(\mathfrak q\) 便为零，所以需要允许 \(C\) 的全部一维特征标参与检测。若对所有 \(\lambda\in\operatorname{Irr}(C)\)，和 \(S_\lambda=\sum_{t\in T}\lambda(t)\) 都属于 \(\mathfrak q\)，取 \(t_0\in T\)，利用循环群的有限几何级数可得
	\[
		\sum_{\lambda\in\operatorname{Irr}(C)} S_\lambda\cdot\lambda(t_0)^{-1} =\sum_{t\in T}\sum_{\lambda\in\operatorname{Irr}(C)} \lambda(tt_0^{-1}) =|C|.
	\]
	这条有限 Fourier 反演等式从各个 \(S_\lambda\) 恢复集合 \(T\) 的示性函数在 \(t_0\) 处的值 \(1\)：内层求和在 \(t=t_0\) 时为 \(|C|\)，其余时为零。左侧模 \(\mathfrak q\) 为零，右侧却因 \(p\mathrel{\vcenter{\hbox{\scalebox{1.25}{\(\nmid\)}}}}|C|\) 非零，矛盾。故至少一个 \(S_\lambda\notin\mathfrak q\)。整数 \(d\) 模 \(\mathfrak q\) 也非零，于是\eqref{b5:eq:elementary-induction-evaluation}不属于 \(\mathfrak q\)。这与该诱导属于 \(I_p(G)_{(p)}\subset\mathfrak m\) 矛盾，故理想等于全环。

	由\eqref{b5:eq:artin-elementary-finite-index}，商 \(R(G)/I_p(G)\) 有限。若其阶被 \(p\) 整除，它含有阶为 \(p\) 的元素，而这样的元素不会在仅反演与 \(p\) 互素整数后消失。这与刚证的局部化商为零矛盾，故指数与 \(p\) 互素。
\end{proof}

平凡特征标的检测确实会在共轭类上失败，下面的二面体群给出一次具体计算。

\begin{proposition}\label{b5:prop:dihedral-fourier-detection}
	设 \(G=D_{10}=\langle r,h\mid r^5=h^2=1,\ hrh=r^{-1}\rangle\) 为十阶二面体群，\(C=\langle r\rangle\)，\(\zeta_{10}=e^{2\pi\mathrm i/10}\)，\(\zeta_5=\zeta_{10}^2\)，\(O=\mathbb Z[\zeta_{10}]\)。设 \(\mathfrak q\subset O_{(2)}\) 为剩余特征是 \(2\) 的极大理想，其中下标 \((2)\) 表示反演所有奇整数。则 \(T=r^G\cap C=\{r,r^{-1}\}\)，且
	\[
		\bigl(\operatorname{Ind}_C^G1_C\bigr)(r)\equiv0\pmod{\mathfrak q},\qquad \bigl(\operatorname{Ind}_C^G\lambda\bigr)(r)\not\equiv0\pmod{\mathfrak q},
	\]
	其中一维复特征标 \(\lambda\) 由 \(\lambda(r)=\zeta_5\) 确定。
\end{proposition}
\begin{proof}
	旋转保持 \(r\)，每个反射把它共轭为 \(r^{-1}\)，所以 \(T=\{r,r^{-1}\}\)，而 \(C_G(r)=C\)。取 \(p=2\) 时，中心化子中的 Sylow \(2\)-子群为 \(1\)，因而局部整诱导证明中的 \(E=C\)、\(d=1\)。\eqref{b5:eq:elementary-induction-evaluation}于是给出
	\[
		\bigl(\operatorname{Ind}_C^G1_C\bigr)(r)=2,\qquad \bigl(\operatorname{Ind}_C^G\lambda\bigr)(r)=\zeta_5+\zeta_5^{-1}.
	\]
	第一个值在剩余特征 \(2\) 中为零。令 \(z\) 为 \(\zeta_5\) 在域 \(O_{(2)}/\mathfrak q\) 中的像。它满足 \(z^5=1\) 及 \(1+z+z^2+z^3+z^4=0\)；若 \(z=1\)，后一式在特征 \(2\) 中变为 \(1=0\)，矛盾。因此 \(z\) 仍有阶 \(5\)。若 \(z+z^{-1}=0\)，则 \(z^2=1\)，又与其阶矛盾，故第二个诱导值不为零。
\end{proof}

这里 \(|T|=2\) 在剩余域中消失，非平凡特征标却仍能检测出 \(T\) 非空。一般证明对全部一维特征标求和，正是为了在某一特征标失效时仍能恢复 \(T\) 的示性函数在每个元素处的取值。

\begin{proposition}\label{b5:prop:brauer-integral-image}
	设 \(G\) 为有限群，\(I(G)\) 为全部初等子群的复表示环诱导像之和。则 \(I(G)=R(G)\)。
\end{proposition}
\begin{proof}
	\(I(G)\) 包含循环诱导，由\cref{b5:cor:artin-finite-index}，加法商群 \(A=R(G)/I(G)\) 有限。对每个素数 \(p\)，它又是 \(R(G)/I_p(G)\) 的商，而\cref{b5:thm:brauer-local-induction}说明后者阶与 \(p\) 互素。故 \(|A|\) 与每个素数互素，只能等于 \(1\)，于是 \(A=0\)。
\end{proof}

这里先用 Artin 得到有限指数，再逐素数排除指数，正是从“某个整数倍可诱导”到“本身可整诱导”的差别。

\begin{corollary}[初等子群上的虚特征标判据]\label{b5:cor:brauer-character-criterion}
	设 \(G\) 为有限群，\(f:G\rightarrow\mathbb C\) 为类函数。则 \(f\in R(G)\)，当且仅当对每个初等子群 \(E\le G\)，限制 \(f|_E\) 属于 \(R(E)\)。
\end{corollary}
\begin{proof}
	正向来自实际表示的限制及可加性。反向中，\cref{b5:prop:brauer-integral-image}把每个不可约 \(\chi\) 写成 \(\sum_j n_j\operatorname{Ind}_{E_j}^G\theta_j\)，其中 \(n_j\in\mathbb Z\)、\(\theta_j\in R(E_j)\)。于是
	\[
		\langle\chi,f\rangle_G =\sum_j n_j\langle\theta_j,f|_{E_j}\rangle_{E_j}\in\mathbb Z.
	\]
	这里使用 Frobenius 互反及两个虚特征标的内积为整数。以不可约特征标为正交规范基展开 \(f\)，所有系数均为整数，故 \(f\in R(G)\)。
\end{proof}

这个论证给出不可约展开系数的整数性。即使把 \(\theta_j\) 选为实际不可约特征标，并且每个 \(f|_{E_j}\) 都是实际特征标，右侧各内积虽然非负，整诱导中的 \(n_j\) 仍可能为负，因而不能由此推出 \(\langle\chi,f\rangle_G\ge0\)。局部限制也确实不能保证这些系数非负。

\begin{proposition}\label{b5:prop:brauer-positive-restriction-failure}
	设 \(G=S_3\)，\(1,\varepsilon,\chi\) 分别为它的平凡、符号和二维不可约复特征标，令 \(f=-1+\varepsilon+\chi\)。则 \(f\) 限制到每个初等子群都是实际复特征标，但 \(f\) 自身不是实际复特征标。
\end{proposition}
\begin{proof}
	\(S_3\) 的真子群阶只能为 \(1,2,3\)，均循环。\(S_3\) 自身不是初等群：若写成 \(P\times C\)，则中心平凡迫使中心中的 \(C\) 为 \(1\)，但阶 \(6\) 又不是素数幂。因此只需检查平凡子群及二阶、三阶子群。取 \(C_2=\langle(12)\rangle\)，令 \(\eta\) 为它的非平凡一维特征标。由 \(\chi\) 在单位元、换位处的值 \(2,0\)，有 \(\chi|_{C_2}=1+\eta\)，而 \(\varepsilon|_{C_2}=\eta\)，故
	\[
		f|_{C_2}=2\eta.
	\]
	取 \(C_3=\langle(123)\rangle\)，令 \(\lambda((123))=\zeta_3\)，其中 \(\zeta_3\) 为本原三次单位根。由三循环处 \(\chi\) 的值为 \(-1\)，有 \(\chi|_{C_3}=\lambda+\lambda^2\)，且 \(\varepsilon|_{C_3}=1\)，所以
	\[
		f|_{C_3}=\lambda+\lambda^2.
	\]
	这两个限制都是实际特征标，共轭子群上的限制由共轭搬运得到同样结论。平凡子群上 \(f(1)=2\)，也为实际特征标。然而 \(\langle1,f\rangle_{S_3}=-1\)，平凡不可约特征标的系数为负，故 \(f\) 不是实际特征标。
\end{proof}

\subsection{初等群表示与一维表示的诱导}
\label{b5:subsec:0504:04}

要将初等群上的不可约表示化为一维诱导，关键在于 \(p\)-群。取\subsecref{b5:subsec:0204:01}中构造的 \(Q_8\) 的二维不可约复表示 \(V\)：在一组基 \(e_1,e_2\) 下，
\[
	\rho(\mathbf i)=\begin{pmatrix}\mathrm i&0\\0&-\mathrm i\end{pmatrix},\qquad
	\rho(\mathbf j)=\begin{pmatrix}0&1\\-1&0\end{pmatrix}.
\]
只限制到中心 \(Z(Q_8)=\{\pm1\}\) 时，全部 \(V\) 都是 \(-1\mapsto-1\) 的同一个权空间，惯性群仍为全群。取 \(x=\mathbf i\)，则 \(xZ(Q_8)\) 位于交换商群 \(Q_8/Z(Q_8)\) 的中心，而
\[
	A=\langle Z(Q_8),\mathbf i\rangle=\langle\mathbf i\rangle\simeq C_4.
\]
这个较大的正规交换子群给出两个不同权：若 \(\lambda(\mathbf i)=\mathrm i\)，则 \(V|_A\cong\lambda\oplus\lambda^{-1}\)，完整的 \(\lambda\)-权空间为 \(V_\lambda=\mathbb Ce_1\)。\(\mathbf j\) 把 \(\mathbf i\) 共轭为其逆，并交换两条权直线，所以 \(I_{Q_8}(\lambda)=A\)。\cref{b5:thm:clifford-correspondence}以 \(N=A\)、\(W=\lambda\) 给出
\[
	V\cong\operatorname{Ind}_A^{Q_8}V_\lambda =\operatorname{Ind}_{\langle\mathbf i\rangle}^{Q_8}\lambda.
\]
在一般 \(p\)-群中，选择比中心大的正规交换子群，也是为了出现不同权，使某个权的惯性群成为真子群，再对它使用阶数归纳；完整权空间此时可以不止一维。
\begin{theorem}\label{b5:thm:p-group-monomial}
	设 \(p\) 为素数，\(P\) 为有限 \(p\)-群。\(P\) 的每个复不可约特征标 \(\theta\) 都可写成 \(\operatorname{Ind}_Q^P\mu\)，其中 \(Q\le P\)，\(\mu\) 为 \(Q\) 的一维复特征标。
\end{theorem}
\begin{proof}
	按 \(|P|\) 归纳。若表示有非平凡核 \(N\)，它来自较小 \(p\)-群 \(P/N\) 的不可约表示。由归纳假设，这个表示等于某个 \(\overline Q\le P/N\) 的一维特征标 \(\overline\mu\) 的诱导。取原像 \(Q\le P\)，令 \(\mu(q)=\overline\mu(qN)\)。映射
	\[
		\mathbb C[P]\otimes_{\mathbb C[Q]}\mathbb C_\mu \longrightarrow \mathbb C[P/N]\otimes_{\mathbb C[\overline Q]}\mathbb C_{\overline\mu}, \qquad g\otimes z\longmapsto(gN)\otimes z
	\]
	与张量平衡关系相容，因为 \(\mu(q)=\overline\mu(qN)\)，并与 \(P\) 作用相容。陪集双射 \(P/Q\simeq(P/N)/\overline Q\) 说明它把一组陪集基双射到另一组，故为同构。于是原表示为 \(\operatorname{Ind}_Q^P\mu\)。

	以下设表示 \(\rho\) 忠实。若 \(P\) 交换，由\cref{b5:cor:schur-algebraically-closed}，全部 \(\rho(g)\) 为标量，不可约性迫使维数为 \(1\)；这也包含归纳起点 \(P=1\)。

	设 \(P\) 非交换。非平凡有限 \(p\)-群中心非平凡，因为类方程中非中心共轭类的大小均被 \(p\) 整除，故中心的阶亦被 \(p\) 整除。将此结论用于非平凡 \(p\)-群 \(P/Z(P)\)，选取 \(xZ(P)\ne1\) 位于其中心。子群 \(A=\bigl\langle Z(P),x\bigr\rangle\) 交换且正规：对任意 \(g\in P\)，有 \(gxg^{-1}\in xZ(P)\)。又因 \(x\notin Z(P)\)，\(A\) 不包含于 \(Z(P)\)。

	由于 \(A\) 交换，它的不可约复表示都是一维特征标。将 \(\rho\) 的表示空间 \(V\) 限制到 \(A\)，按这些特征标的完整等型分量分解：
	\[
		V=\bigoplus_\lambda V_\lambda,\qquad V_\lambda=\{\,v\in V:\rho(a)v=\lambda(a)v\;\text{对所有}\;a\in A\,\}.
	\]
	每个非零 \(V_\lambda\) 中可以有多个同类型的一维成分，因而 \(\dim V_\lambda\) 不必为一。若只出现一个特征标，则每个 \(a\in A\) 都作用为标量 \(\rho(a)=\lambda(a)I\)。对任意 \(g\in P\)，于是有
	\[
		\rho([a,g]) =\rho(a)\rho(g)\rho(a)^{-1}\rho(g)^{-1}=I.
	\]
	忠实性给出 \([a,g]=1\)，故 \(A\subseteq Z(P)\)，矛盾。因此至少出现两个权。

	取其中一个 \(\lambda\)，其惯性群 \(J=I_P(\lambda)\) 为真子群；这是因为\cref{b5:thm:clifford-restriction}断言不可约表示中的所有权构成一个 \(P\)-轨道，而这里已有至少两个权。完整权空间 \(V_\lambda\) 被 \(J\) 保持：对 \(j\in J\)、\(a\in A\)、\(v\in V_\lambda\)，有
	\[
		\rho(a)\rho(j)v =\rho(j)\rho(j^{-1}aj)v =\lambda(j^{-1}aj)\rho(j)v =\lambda(a)\rho(j)v.
	\]
	它作为 \(J\)-表示也不可约。事实上，若 \(0\ne Z\subseteq V_\lambda\) 是 \(J\)-子模，则对 \(P/J\) 的一组代表 \(g\)，各 \(\rho(g)Z\) 位于不同的权空间；它们的直和被 \(P\) 保持，故由 \(V\) 的不可约性等于 \(V\)。取其中的 \(\lambda\)-权分量便得 \(Z=V_\lambda\)。因此\cref{b5:thm:clifford-correspondence}在这里给出的 \(U\) 正是整个 \(V_\lambda\)，且 \(\rho\simeq\operatorname{Ind}_J^P U\)。对较小的 \(p\)-群 \(J\) 应用归纳假设，得 \(U\simeq\operatorname{Ind}_Q^J\mu\)。\cref{b5:thm:induction-transitivity}遂给出 \(\rho\simeq\operatorname{Ind}_Q^P\mu\)。
\end{proof}

\begin{proposition}\label{b5:prop:elementary-monomial}
	有限初等群 \(E\) 的每个复不可约特征标都是其某个初等子群的一维复特征标的诱导。
\end{proposition}
\begin{proof}
	写 \(E=P\times C\)。在不可约表示中，中心子群 \(C\) 作用为某个一维特征标 \(\lambda\)，而限制到 \(P\) 是不可约表示 \(\theta\)：若有真 \(P\)-不变子空间，标量的 \(C\)-作用也保持它，与 \(E\)-不可约性矛盾。由\cref{b5:thm:p-group-monomial}，\(\theta=\operatorname{Ind}_Q^P\mu\)。取 \(P/Q\) 的陪集基，便有
	\[
		\theta\otimes\lambda =\operatorname{Ind}_{Q\times C}^{P\times C}(\mu\otimes\lambda).
	\]
	其中 \(Q\times C\) 仍是 \(p\)-初等群，且 \(\mu\otimes\lambda\) 是它的一维特征标，得到所需表达。
\end{proof}

\begin{proof}[\cref{b5:thm:brauer-induction}的证明]
	由\cref{b5:prop:brauer-integral-image}，任意 \(\alpha\in R(G)\) 是初等子群虚特征标诱导的整数线性组合；把这些虚特征标展开为不可约特征标的整数线性组合，便得到定理的第一项结论。对每个不可约初等子群特征标应用\cref{b5:prop:elementary-monomial}，将其写成同一初等群内某个初等子群的一维诱导。最后用\cref{b5:thm:induction-transitivity}把两次诱导合成一次，即得全部系数仍为整数的一维表述。
\end{proof}

\subsection{负系数与虚特征标的例子}
\label{b5:subsec:0504:05}

取 \(G=S_3\)，其不可约特征标为 \(1,\varepsilon,\chi\)，次数分别为 \(1,1,2\)。设 \(C_2=\bigl\langle(12)\bigr\rangle\)，\(C_3=\bigl\langle(123)\bigr\rangle\)，取 \(C_3\) 的非平凡一维特征标 \(\lambda\)。由限制及 Frobenius 互反，
\[
	\operatorname{Ind}_{C_2}^{S_3}1=1+\chi,\qquad \operatorname{Ind}_{C_3}^{S_3}\lambda=\chi.
\]
第一式也可从三点置换表示中分离常数向量得到；第二式在三类上的值依次为 \(2,0,-1\)。所以
\begin{equation}\label{b5:eq:s3-brauer-negative}
	1=\operatorname{Ind}_{C_2}^{S_3}1 -\operatorname{Ind}_{C_3}^{S_3}\lambda.
\end{equation}
两个子群都循环，因而都是允许的初等子群。此式移项后表示一个三维实际表示同构于一维平凡表示与一个二维表示的直和；负号没有制造负维数的表示空间。

\ProseParagraphBreak
\(S_3\) 自身不是初等群，所有初等子群都是真子群，故每个一维诱导的次数至少为 \(2\)。因此 \(1\) 不可能是这些诱导的非负整数线性组合。这也说明负系数不是上述公式偶然写得不巧，而是整诱导定理必须允许的现象。

% !TeX root = ../../Book5.tex
% 纲要标题：### 5.5 共轭类、特征标约束与 Burnside 可解性定理
% 纲要位置：工作记录/计划纲要.md，第 3882 行。
% 纲要细目（写作范围）：
% * 先准确陈述 Burnside 的 \(p^aq^b\) 阶群可解性定理，再完成下列共轭类论证；允许其中一个指数为零，并分别处理平凡群和素数幂阶群。
% * 从类和作用和代数整数性建立群论所需的特征标消失或标量作用引理；注明次数与共轭类大小互素的使用位置
% * 证明有限非交换单群不存在非平凡的素数幂大小共轭类，并处理类大小为1的中心情形
% * 对 \(|G|=p^aq^b\) 利用 Sylow 子群中心中的元素构造所需共轭类，再按群阶归纳并使用可解性的扩张稳定性完成证明
% * 以中心化子指数说明素数幂共轭类对有限非交换单群的限制

\section{共轭类、特征标约束与 Burnside 可解性定理}
\label{b5:sec:0505}

整性还能把“某个共轭类很小”变成正规子群的存在性。以下从特征标正交、\cref{b5:prop:central-character-integrality}与代数整数的基本运算出发。

\subsection{\texorpdfstring{Burnside 的 \(p^aq^b\) 阶群可解性定理}{Burnside 的 paqb 阶群可解性定理}}
\label{b5:subsec:0505:01}

对群 \(G\)，记 \(G'=\langle xyx^{-1}y^{-1}:x,y\in G\rangle\)，递归定义 \(G^{(0)}=G\)、\(G^{(n+1)}=\bigl(G^{(n)}\bigr)'\)。若某个 \(G^{(r)}=1\)，则称 \(G\) 为\term[kejiequn]{可解群}[solvable group]。这一定义等价于存在每个相邻商群均交换的有限次正规列
\[
	G=G_0\supseteq G_1\supseteq\cdots\supseteq G_r=1, \qquad G_{i+1}\trianglelefteq G_i.
\]
导出列本身给出一个方向；反过来，\(G_i/G_{i+1}\) 交换意味着 \(G_i'\subseteq G_{i+1}\)，由归纳得到 \(G^{(i)}\subseteq G_i\)，因而 \(G^{(r)}=1\)。

\begin{theorem}[Burnside 可解性定理]\label{b5:thm:burnside-solvability}
	设 \(p,q\) 为素数，\(a,b\) 为非负整数。每个阶为 \(p^a q^b\) 的有限群 \(G\) 都是可解群。
\end{theorem}
此结论称为\term[Burnsidekejie]{Burnside 可解性定理}[Burnside's \(p^a q^b\) theorem]。允许 \(a=0\)、\(b=0\)；若 \(p=q\)，结论仍归入素数幂阶情形。证明在本节先建立共轭类工具，再于\subsecref{b5:subsec:0505:04}完成。

\ProseParagraphBreak
定理并不说 \(G\) 交换，也不说全部 Sylow 子群正规。例如 \(S_3\) 阶为 \(2\cdot3\)，其导出群为 \(A_3\)，故可解，但三阶以外的 Sylow 子群并不正规。这里的目标是在导出列中逐步消去非交换部分。

\subsection{类和作用与特征标消失引理}
\label{b5:subsec:0505:02}

\begin{lemma}[单位根平均值]\label{b5:lem:root-unity-average}
	设 \(d\ge1\) 为整数，\(\xi_1,\ldots,\xi_d\) 为复单位根。若
	\[
		\alpha=\frac{\xi_1+\cdots+\xi_d}{d}
	\]
	是代数整数，则或者 \(\alpha=0\)，或者 \(\xi_1=\cdots=\xi_d\)。
\end{lemma}
\begin{proof}
	若这些单位根不全相同，三角不等式的等号条件给出 \(|\alpha|<1\)。选取公倍数 \(m\)，使各 \(\xi_i\) 都属于 \(\mathbb Q(\zeta_m)\)。\(\alpha\) 的每个代数共轭仍是 \(d\) 个单位根的平均，故绝对值不超过 \(1\)。这里的共轭描述可这样核验：\(\mathbb Q(\alpha)\) 到 \(\mathbb C\) 的任一嵌入都可延拓到 \(\mathbb Q(\alpha,\zeta_m)\)，设该嵌入为 \(\sigma\)，\(f\) 为 \(\zeta_m\) 在 \(\mathbb Q(\alpha)\) 上的最小多项式。对 \(f\) 的系数施行 \(\sigma\) 得到 \(\sigma(f)\)，再选一个复根 \(\beta\)。映射 \(\sum a_iT^i\mapsto\sum\sigma(a_i)\beta^i\) 通过商域 \(\mathbb Q(\alpha)[T]/(f)\simeq\mathbb Q(\alpha,\zeta_m)\) 分解；它保幺，因定义域是域而单射，故给出所需延拓。延拓保持 \(\zeta_m^m=1\)，因而保持“是单位根”这一性质。

	若 \(\alpha\ne0\)，其全部共轭根的乘积非零；由\cref{b5:lem:algebraic-integer-arithmetic}该乘积为整数。但其中一个因子的绝对值严格小于 \(1\)，其余不超过 \(1\)，故乘积的绝对值严格小于 \(1\)，矛盾。
\end{proof}

证明中的乘积估计同时使用了一个共轭的绝对值严格小于 \(1\) 和全部共轭的绝对值不超过 \(1\)。只控制一个共轭并不足够：\(\beta=2-\sqrt3\) 满足首一整系数方程 \(\beta^2-4\beta+1=0\)，是绝对值小于 \(1\) 的非零代数整数；它的另一个共轭 \(2+\sqrt3\) 却大于 \(1\)，所以两个共轭的乘积仍为整数 \(1\)。单位根平均值的特殊性在于，每个嵌入都把各单位根送到单位圆上，从而使全部共轭同时受到模长限制。

\begin{lemma}[消失或标量作用]\label{b5:lem:coprime-class-scalar}
	设 \(G\) 为有限群，\(\rho:G\rightarrow\operatorname{GL}_{\mathbb C}(V)\) 为次数 \(d\) 的复不可约表示，特征标为 \(\chi\)，\(g\in G\)。设 \(K=g^G\) 且 \(\gcd(d,|K|)=1\)。则或者 \(\chi(g)=0\)，或者 \(\rho(g)\) 为标量矩阵。
\end{lemma}
\begin{proof}
	取整数 \(a,b\)，使 \(a|K|+bd=1\)。由\cref{b5:prop:central-character-integrality}，\(|K|\cdot\chi(g)/d\) 是代数整数；\(\chi(g)\) 也由\cref{b5:prop:character-values-integral}为代数整数。因此
	\[
		\frac{\chi(g)}d =a\cdot\frac{|K|\cdot\chi(g)}d+b\cdot\chi(g)
	\]
	是代数整数。这正是互素条件的使用处。\(\rho(g)\) 可对角化，其特征值 \(\xi_1,\ldots,\xi_d\) 为单位根；由\cref{b5:lem:root-unity-average}，或者它们的和为零，或者全部相同。后者连同可对角化性给出标量矩阵。
\end{proof}

不能删去互素条件。例如 \(S_3\) 的二维不可约表示在三循环处取值 \(-1\)，并非零；对应矩阵的两个特征值为互不相同的本原三次单位根，所以不是标量。该共轭类大小为 \(2\)，恰与表示次数有公因子。

\subsection{非交换单群的素数幂大小共轭类}
\label{b5:subsec:0505:03}

\begin{proposition}\label{b5:prop:prime-power-class-normal-subgroup}
	设 \(G\) 为有限群，某个共轭类 \(K\) 的大小为 \(p^r\)，其中 \(p\) 为素数、\(r\ge1\)。则 \(G\) 有非平凡真正规子群。特别地，有限非交换单群没有这样的共轭类。
\end{proposition}
\begin{proof}
	取 \(g\in K\)，则 \(g\ne1\)。由\cref{b5:thm:regular-decomposition}，正则特征标在非单位元处为零，故
	\begin{equation}\label{b5:eq:burnside-regular-sum}
		0=\sum_{\chi\in\operatorname{Irr}(G)}\chi(1)\chi(g).
	\end{equation}
	我们先证明存在非平凡不可约 \(\chi\)，满足 \(p\nmid\chi(1)\) 且 \(\chi(g)\ne0\)。否则，\eqref{b5:eq:burnside-regular-sum}除去平凡特征标的贡献 \(1\)，其余可能非零的项的次数均被 \(p\) 整除，从而
	\[
		-\frac1p =\sum_{\substack{\chi\in\operatorname{Irr}(G)\\p\mid\chi(1)}} \frac{\chi(1)}p\,\chi(g)
	\]
	为代数整数。这与\cref{b5:lem:algebraic-integer-arithmetic}的有理数判别矛盾。

	取这样的 \(\chi\) 及表示 \(\rho\)。因为 \(\gcd\bigl(\chi(1),p^r\bigr)=1\) 且 \(\chi(g)\ne0\)，\cref{b5:lem:coprime-class-scalar}说明 \(\rho(g)\) 为标量。所有 \(K\) 中的元素共轭于 \(g\)，所以它们在 \(\rho\) 下给出同一个标量 \(cI\)，其中 \(c\ne0\)。这个标量未必为 \(1\)，故 \(K\) 本身未必包含于 \(\ker\rho\)，直接取 \(\langle K\rangle\) 也不能保证得到真子群。为消去共同标量，考虑类中两个元素的商：
	\[
		\rho(xy^{-1})=\rho(x)\rho(y)^{-1}=cI(cI)^{-1}=I, \qquad x,y\in K.
	\]
	因此令
	\[
		N=\langle xy^{-1}:x,y\in K\rangle.
	\]
	共轭保持 \(K\)，故 \(N\triangleleft G\)；又 \(N\subset\ker\rho\)。由于 \(\rho\) 非平凡，\(\ker\rho\ne G\)，所以 \(N\ne G\)。而 \(|K|>1\) 给出不同的 \(x,y\)，此时 \(xy^{-1}\ne1\)，所以 \(N\ne1\)。
\end{proof}

大小为 \(1\) 的类须另行处理，它对应中心元素。有限非交换单群的中心只能为 \(1\)：中心若非平凡，单性使整个群交换；有限交换单群则必为素数阶循环群。因此非交换单群的非单位元既不能有大小为 \(1\) 的类，也不能有更大的素数幂大小的类。

\subsection{Sylow 中心、群阶归纳与可解扩张}
\label{b5:subsec:0505:04}

先记录可解性的群论运算，避免把归纳的最后一步藏在一句“因此可解”中。
\begin{lemma}\label{b5:lem:solvable-extensions}
	设 \(G\) 为群，\(N\triangleleft G\)。若 \(N\) 与 \(G/N\) 都可解，则 \(G\) 可解。可解群的子群和商群也可解。对每个素数 \(p\)，每个有限 \(p\)-群都可解。
\end{lemma}
\begin{proof}
	若 \((G/N)^{(r)}=1\)，商映射保持交换子，故 \(G^{(r)}\subset N\)。若 \(N^{(s)}=1\)，则 \(G^{(r+s)}\subset N^{(s)}=1\)。子群 \(H\le G\) 满足 \(H^{(i)}\subset G^{(i)}\)，而满同态把导出列映为商群的导出列，由此得到另外两项稳定性。

	对有限 \(p\)-群 \(P\) 按阶归纳。若 \(P\ne1\)，其类方程中非中心类的大小均被 \(p\) 整除，故 \(|Z(P)|\) 被 \(p\) 整除，中心非平凡。若 \(Z(P)=P\)，群交换；否则 \(P/Z(P)\) 阶较小，由归纳假设可解。中心自身交换，再由已经证明的扩张性质得 \(P\) 可解。
\end{proof}

\begin{proof}[\cref{b5:thm:burnside-solvability}的证明]
	平凡群可解；若 \(a=0\)、\(b=0\) 或 \(p=q\)，由\cref{b5:lem:solvable-extensions}的有限 \(p\)-群结论可得。以下设 \(p\ne q\) 且 \(a,b>0\)，按 \(|G|\) 归纳。

	若 \(G\) 有非平凡真正规子群 \(N\)，则 \(N\) 与 \(G/N\) 的阶仍只含素因子 \(p,q\)，且都小于 \(|G|\)。归纳假设使两者可解，再由\cref{b5:lem:solvable-extensions}得 \(G\) 可解。只需排除 \(G\) 为单群的情形。它不可能是交换单群，因为交换单群只有一个素因子。

	取 Sylow \(p\)-子群 \(P\)。由 \(p\)-群的类方程，选 \(1\ne x\in Z(P)\)。于是 \(P\subset C_G(x)\)。由于 \(|P|=p^a\)，Lagrange 定理给出 \(p^a\bigm\vert |C_G(x)|\)，而 \(|C_G(x)|\) 又整除 \(p^aq^b\)，所以对某个 \(0\le c\le b\) 有 \(|C_G(x)|=p^aq^c\)。因此
	\[
		|x^G|=\bigl[G:C_G(x)\bigr] =\frac{p^aq^b}{p^aq^c}=q^{b-c}.
	\]
	选择 Sylow 子群中心中的元素，恰使共轭类大小中的整个 \(p\)-部分被消去；这只需要 \(x\) 与 \(P\) 中每个元素交换，不要求 \(P\) 正规。若这个指数为 \(1\)，则 \(x\in Z(G)\setminus\{1\}\)，与非交换单性矛盾；若指数大于 \(1\)，则\cref{b5:prop:prime-power-class-normal-subgroup}给出非平凡真正规子群，同样矛盾。因此 \(G\) 不可能为单群，归纳即得可解性。
\end{proof}

在 \(S_3\) 中取 \(p=2,q=3\)、\(P=\langle(12)\rangle\)。这个 Sylow 子群不正规，但 \(Z(P)=P\)，而 \(x=(12)\) 的中心化子正是 \(P\)，故 \(K=x^{S_3}=\{(12),(13),(23)\}\) 的大小为 \(3\)。一维符号表示 \(\varepsilon\) 把每个换位都送到 \(-1\)，所以 \(K\) 不包含于它的核；两个换位的商却被送到 \(1\)。例如 \((12)(13)^{-1}=(132)\)，而所有这些商都是偶置换，因而
\[
	N=\langle xy^{-1}:x,y\in K\rangle=A_3=\ker\varepsilon.
\]
于是得到正规列 \(1\triangleleft A_3\triangleleft S_3\)，其中 \(A_3\simeq C_3\)、\(S_3/A_3\simeq C_2\)。两个群均交换，\cref{b5:lem:solvable-extensions}便给出 \(S_3\) 可解。这个例子把非正规 Sylow 子群的中心、素数幂大小的共轭类、消去共同标量和可解扩张联系在同一个计算中。

\ProseParagraphBreak
整性与共轭类步骤可对照 Etingof 等\cite[Theorems 5.4.3, 5.4.4, 5.4.6; Lemmas 5.4.5, 5.4.7, pp. 98--100]{EtingofEtAl2011RepresentationTheory}；这里的最后一步以 Sylow 中心显式构造所需共轭类。

\subsection{素数幂共轭类对单群的限制}
\label{b5:subsec:0505:05}

有限非交换单群的每个非单位元 \(g\) 都必须满足：中心化子指数 \([G:C_G(g)]\) 既不为 \(1\)，也不是任何素数的正整数次幂。因此每个非平凡共轭类的大小至少含两个不同素因子。这个限制施加在共轭类大小上，而不只是群阶上；对只含两个素因子的群，Sylow 子群中心提供的元素使中心化子指数只剩一个素因子，正是在这一点与单性矛盾。下一节将用可解性进一步确定合成因子与极小正规子群的形状。

% !TeX root = ../../Book5.tex
% 纲要标题：### 5.6 群论应用与三种 Burnside 定理的比较
% 纲要位置：工作记录/计划纲要.md，第 3890 行。
% 纲要细目（写作范围）：
% * Burnside 的 p^a q^b 阶群可解性定理
% * 与第二卷 Burnside 矩阵代数定理的区别
% * 结构结论如何由特征标约束推出
% * 本节运用第5.5节已经证明的 \(p^aq^b\) 阶群可解性定理，整理群论推论及适用范围；将其与轨道计数引理、矩阵代数定理逐项比较
% ---

\section{群论应用与三种 Burnside 定理的比较}
\label{b5:sec:0506}

可解性是一个结构结论，尚未给出具体群的分类。使用它时，通常再取极小或极大正规子群，把群拆成较简单的扩张。与此同时，特征标方法还可以研究置换作用和矩阵代数；应当分清每次得到了什么。

\subsection{\texorpdfstring{\(p^aq^b\) 阶群可解性定理的应用}{paqb 阶群可解性定理的应用}}
\label{b5:subsec:0506:01}

\begin{corollary}\label{b5:cor:two-prime-composition}
	设 \(p,q\) 为素数，\(a,b\ge0\) 为整数，\(G\) 为阶大于 \(1\) 且 \(|G|=p^a q^b\) 的有限群。则 \(G\) 有素数阶循环商群，并且其每个合成因子均同构于 \(C_p\) 或 \(C_q\)。特别地，有限非交换单群的阶至少含三个不同素因子。
\end{corollary}
\begin{proof}
	由\cref{b5:thm:burnside-solvability}，\(G\) 可解。取极大真正规子群 \(N\)，则 \(G/N\) 为非平凡单群，并由\cref{b5:lem:solvable-extensions}仍可解。非平凡单可解群必须交换：其导出群若为自身，导出列不会终止；否则单性使导出群为 \(1\)。有限交换单群为素数阶循环群，阶又整除 \(|G|\)，故只能为 \(p\) 或 \(q\)。任取一条合成列，各合成因子 \(G_i/G_{i+1}\) 既单又可解，因为子群与商群保持可解性；其阶也整除 \(|G|\)。上述论证遂使每个因子都同构于 \(C_p\) 或 \(C_q\)。最后一项是\cref{b5:thm:burnside-solvability}的直接逆否应用。
\end{proof}

这并不使扩张自动分裂，也不能由合成因子的清单恢复乘法。例如 \(C_4\) 与 \(C_2\times C_2\) 都有两个 \(C_2\) 合成因子，却不同构。表示论给出的限制常要与具体正规子群及共轭作用一起使用。

\subsection{与 Burnside 矩阵代数定理的比较}
\label{b5:subsec:0506:02}

\cref{b5:thm:complex-group-algebra}已经给出复群代数的矩阵块分解
\[
	\Phi:\mathbb C[G]\xrightarrow{\sim}\prod_i\operatorname{End}_{\mathbb C}(S_i).
\]
把 \(\Phi\) 与第 \(i\) 个因子的投影复合，就是不可约表示 \(S_i\) 上的群代数作用 \(\rho_i\)，因此这个作用是满射。群代数元素是群元素的复线性组合，所以
\[
	\rho_i\bigl(\mathbb C[G]\bigr) =\operatorname{span}_{\mathbb C}\{\rho_i(g):g\in G\} =\operatorname{End}_{\mathbb C}(S_i).
\]
下面的矩阵代数定理把这一现象推广到代数闭域上任意不可约作用的含幺矩阵子代数，所讨论的代数无需来自有限群。
\begin{theorem}[Burnside 矩阵代数定理]\label{b5:thm:burnside-matrix-algebra}
	设 \(k\) 为代数闭域，\(V\ne0\) 为有限维 \(k\)-向量空间，\(A\subset\operatorname{End}_k(V)\) 为含恒等映射的 \(k\)-子代数。若 \(V\) 没有非零真 \(A\)-不变子空间，则 \(A=\operatorname{End}_k(V)\)。
\end{theorem}
\term[Burnsidejuzhendaishu]{Burnside 矩阵代数定理}[Burnside's matrix algebra theorem]的假设中没有有限群，也没有两个素因子。
\begin{proof}
	首先，每个与 \(A\) 交换的自同态 \(T\) 都是标量。由于 \(k\) 代数闭，\(T\) 有特征值 \(\lambda\)，非零子空间 \(\ker(T-\lambda I)\) 为 \(A\)-不变，故等于 \(V\)。

	对一个非零向量 \(v\)，不可约性已给出 \(Av=V\)，但这只说明能指定这一个向量的像。要得到全部自同态，必须用同一个 \(a\in A\) 同时指定一组基的所有像。取 \(V\) 的基 \(v_1,\ldots,v_d\)，这些能同时实现的像组成
	\[
		M=\bigl\{(av_1,\ldots,av_d):a\in A\bigr\}\subset V^d.
	\]
	它是 \(A\)-子模；所需结论就是 \(M=V^d\)。

	为排除 \(M\) 是真子模，所需的工具是一个把它送到零的非零 \(A\)-线性映射。先看 \(V^2\) 中的对角子模 \(W=\{(v,v):v\in V\}\)。它不是任何一个坐标副本，但
	\[
		(v,w)=(w,w)+(v-w,0)
	\]
	说明 \(V^2=W\oplus(V\oplus0)\)：两个子模的交为零，且上述分解覆盖全部向量。沿 \(W\) 投到这个补空间，再取第一坐标，便得到非零 \(A\)-线性映射 \(\ell(v,w)=v-w\)，满足 \(\ell(W)=0\)。

	一般地，\(A\) 在 \(V^r\) 上逐坐标作用，每个坐标副本都是同构于 \(V\) 的简单 \(A\)-模。对任意子模 \(W\)，从这些坐标模中取尽可能多的项，使其直和 \(U\) 满足 \(U\cap W=0\)。若 \(U+W\ne V^r\)，便有一个坐标简单模 \(S\) 不包含于 \(U+W\)；由简单性，\(S\cap(U+W)=0\)。若 \(u+s\in W\)，其中 \(u\in U,s\in S\)，则 \(s\in S\cap(U+W)\)，故 \(s=0\)，再由 \(U\cap W=0\) 得 \(u=0\)。因此 \((U\oplus S)\cap W=0\)，可以继续加入 \(S\)，与所取项数最大矛盾。故 \(V^r=W\oplus U\)。若 \(W\ne V^r\)，则 \(U\ne0\)。因为 \(W,U\) 都是 \(A\)-子模，沿 \(W\) 到 \(U\) 的投影 \(\pi_U:V^r\rightarrow U\) 是 \(A\)-线性的；再投到 \(U\) 的一个坐标简单直和项 \(S\simeq V\)，复合
	\[
		V^r\xrightarrow{\pi_U}U\longrightarrow S\xrightarrow{\sim}V
	\]
	为非零 \(A\)-线性映射，且在 \(W\) 上为零。

	若 \(M\ne V^d\)，前一事实给出零化 \(M\) 的非零 \(A\)-线性映射 \(\ell:V^d\rightarrow V\)。把 \(\ell\) 限制到第 \(i\) 个坐标副本，得到 \(V\) 的 \(A\)-自同态，由第一段必为标量 \(c_iI\)，故 \(\ell(w_1,\ldots,w_d)=\sum_i c_iw_i\)，其中 \(c_i\) 不全为零。由于 \(I\in A\)，有 \((v_1,\ldots,v_d)\in M\)；而 \(\ell(M)=0\) 给出 \(\sum_i c_i v_i=0\)，与基的线性无关性矛盾。因此 \(M=V^d\)。任意指定各 \(v_i\) 的像都能由某个 \(a\in A\) 实现，故 \(A\) 是全部自同态代数。
\end{proof}

对复不可约有限群表示 \(\rho\)，代数 \(A=\operatorname{span}_{\mathbb C}\rho(G)\) 含有 \(I=\rho(1)\)，且由 \(\rho(g)\rho(h)=\rho(gh)\) 对乘法封闭；它的不变子空间恰是 \(G\)-不变子空间。因此矩阵代数定理也给出本小节开始时从群代数分解得到的结论。例如 \(S_3\) 的二维标准表示的群像只有六个矩阵，但它们的复线性组合生成全部四维空间 \(M_2(\mathbb C)\)。所得到的是线性生成空间等于全部矩阵代数，群像本身仍只是一个有限群。

\ProseParagraphBreak
若底域不是代数闭域，结论可能失败：\(\mathbb C\) 通过实线性乘法作用在 \(\mathbb R^2\) 上不可约，但所得二维实子代数不等于四维的 \(M_2(\mathbb R)\)。

\subsection{由特征标约束推出群结构}
\label{b5:subsec:0506:03}

对非平凡正规子群 \(N\triangleleft G\)，极小正规性是指 \(N\) 内没有非平凡真子群在整个 \(G\) 中正规。它约束的是哪些子群被全部 \(G\)-共轭保持，并不要求 \(N\) 自身为单群。

\begin{proposition}\label{b5:prop:minimal-normal-solvable}
	设 \(G\) 为有限可解群，\(N\triangleleft G\) 为非平凡极小正规子群。则存在素数 \(\ell\) 与正整数 \(r\)，使 \(N\simeq C_\ell^r\)。把群乘法视作加法、整数次幂视作 \(\mathbb F_\ell\) 的标量作用，选基后共轭给出不可约的线性表示 \(G\rightarrow\operatorname{GL}_r(\mathbb F_\ell)\)。若另有素数 \(p,q\) 及非负整数 \(a,b\) 满足 \(|G|=p^a q^b\)，则 \(\ell\) 为 \(p\) 或 \(q\)。更一般地，对任意有限群 \(G\)、素数 \(\ell\) 与正整数 \(r\)，若 \(N\triangleleft G\) 且 \(N\simeq C_\ell^r\)，则 \(N\) 的极小正规性等价于上述共轭作用不可约。
\end{proposition}
\begin{proof}
	由\cref{b5:lem:solvable-extensions}，\(N\) 可解。其导出群 \(N'\) 是 \(N\) 的特征子群，因而在 \(G\) 中正规。极小性说明 \(N'=1\) 或 \(N'=N\)；后一情形使非平凡的 \(N\) 的导出列永不终止，故 \(N\) 交换。选一个整除 \(|N|\) 的素数 \(\ell\)。有限交换群中所有 \(\ell\)-幂阶元素组成非平凡特征子群，极小性使其等于 \(N\)，所以 \(N\) 为交换 \(\ell\)-群。其子群 \(N[\ell]=\{x:x^\ell=1\}\) 同样非平凡且特征，故等于 \(N\)。用群乘法作加法、用整数次幂作 \(\mathbb F_\ell\) 的标量作用，\(N\) 成为有限维 \(\mathbb F_\ell\)-向量空间，因而 \(N\simeq C_\ell^r\)。若 \(|G|=p^aq^b\)，由 \(|N|\bigm\vert |G|\) 得到 \(\ell\in\{p,q\}\)。

	每个群自同构保持群乘法和整数次幂，因此共轭 \(\gamma(g)(n)=gng^{-1}\) 是 \(N\) 的 \(\mathbb F_\ell\)-线性自同构，且 \(\gamma(gh)=\gamma(g)\gamma(h)\)，所以确实给出线性表示。\(N\) 的每个子群都对整数次幂封闭，因而正是这个向量空间的子空间；其中被共轭作用保持的子空间，恰好是包含于 \(N\) 的 \(G\)-正规子群。极小正规性使它们只能为零子空间或 \(N\)，故作用不可约。反过来，若共轭作用不可约，这一对应也说明 \(N\) 内没有非平凡真 \(G\)-正规子群，故 \(N\) 极小正规。此对应只用到了 \(N\simeq C_\ell^r\) 和 \(N\triangleleft G\)，给出最后的一般等价。
\end{proof}

以 \(A_4\) 的正规子群
\[
	N=V_4=\{1,(12)(34),(13)(24),(14)(23)\}\simeq C_2^2
\]
为例。三个双换位都是二阶元素，两个不同双换位的积是第三个，而共轭保持双换位的循环型，所以 \(V_4\) 确为 \(A_4\) 的正规子群。它的三个二阶子群在交换群 \(V_4\) 内均正规，所以 \(V_4\) 自身不是单群。但 \((123)\) 的共轭作用依次把 \((12)(34)\) 送到 \((14)(23)\)，再送到 \((13)(24)\)，再送回 \((12)(34)\)。三个二阶子群因此相互轮换，没有一个在 \(A_4\) 中正规；而它们又是 \(V_4\) 的全部非平凡真子群，所以 \(V_4\) 是 \(A_4\) 的极小正规子群。把 \(V_4\) 看成 \(\mathbb F_2^2\)，这三个非单位元就是三个非零向量，每条一维子空间只含其中一个；共轭作用不保持任何一维子空间，故这个二维表示不可约。这正说明极小正规子群可以有 \(r>1\)。

\subsection{群论推论、适用范围与同名定理}
\label{b5:subsec:0506:04}

轨道计数公式也可完全放在特征标语言中。若有限群 \(G\) 作用在有限集合 \(X\) 上，置换特征标为 \(\pi_X(g)=|X^g|\)。不变向量在各轨道上系数恒定，因此其维数为轨道数；也可直接计算
\[
	\dim\mathbb C[X]^G =\langle1_G,\pi_X\rangle_G =\frac1{|G|}\sum_{g\in G}|X^g|.
\]
从双重计数看，每个轨道的贡献是 \(\sum_{x\in\mathcal O}|G_x|=|\mathcal O|\cdot|G_x|=|G|\)，同样得到公式。

\ProseParagraphBreak
三种结论于是各有明确用途：轨道计数求不变向量的维数或轨道数；矩阵代数定理识别不可约作用所张成的全部代数；\(p^a q^b\) 定理由类和整性推出群的可解性。前两者对 \(A_5\) 一样适用，但 \(|A_5|=60=2^2\cdot3\cdot5\) 已超出第三者的两个素因子假设。对阶含三个以上素因子的群，仅凭本章定理既不能断言可解，也不能断言不可解。


\begin{chaptersummary}
	表示环是以不可约特征标为基的整数格，张量积给出乘法。特征标值由单位根相加而成，类和标量也满足首一整系数方程；将这些整性与有理数判别结合，便得到不可约次数整除群阶。

	\ProseParagraphBreak
	Artin 诱导由循环子群上的限制检测零类函数，再利用整数矩阵的秩获得有理张成。它允许一般的一维特征标；若只从平凡特征标诱导，则必须相应限制为在有理共轭类上常值的函数。循环诱导格的有限指数可能大于 \(1\)，不能把清分母误当成整系数结论。

	\ProseParagraphBreak
	Brauer 诱导扩大到初等子群。对每个素数 \(p\)，局部化表示环 \(R(G)_{(p)}\) 的极大理想都可通过将特征标在某个 \(p\)-正则元素处的取值模去一个剩余特征为 \(p\) 的极大理想来检测；中心化子中的 Sylow 子群与该元素生成的循环群组成所需初等子群。有限 Fourier 求和排除了 \(I_p(G)_{(p)}\) 落入极大理想的可能，得到 \(I_p(G)_{(p)}=R(G)_{(p)}\)，再由逐素数检测回到整数。Clifford 对应给出的 \(p\)-群单项式性则将初等子群的特征标进一步写成一维特征标的诱导。

	\ProseParagraphBreak
	Burnside 可解性定理沿独立的论证使用整性：次数与类大小互素时，特征标或消失，或对应矩阵为标量；正则特征标之和随后从素数幂大小的类中检出正规子群。Sylow 中心和可解扩张完成按群阶的归纳。其结论、轨道计数公式和不可约矩阵代数定理分别处理群结构、群作用计数和线性算子代数，不宜只凭同名而互相替代。
\end{chaptersummary}

\BookExercises

\begin{exercise}\label{b5:ex:05-cyclic-ring}
	设 \(C_n=\langle c\rangle\)，一维特征标 \(\lambda\) 满足 \(\lambda(c)=e^{2\pi i/n}\)。证明 \(R(C_n)\simeq\mathbb Z[T]/(T^n-1)\)，并在此描述次数映射及其核。
\end{exercise}
\begin{solution}
	循环群的不可约特征标为 \(1,\lambda,\ldots,\lambda^{n-1}\)，乘法使指数模 \(n\) 相加。因此 \(T\mapsto\lambda\) 给出满同态。商环中的每个元素唯一写成次数小于 \(n\) 的整系数多项式，而这些特征标在加法群中线性无关，所以同态单射。次数把每个 \(\lambda^j\) 送到 \(1\)，对应 \(f(T)\mapsto f(1)\)。整系数除法给出 \(f(1)=0\) 当且仅当 \(f\) 被 \(T-1\) 整除，故核由 \(T-1\) 的剩余类生成。
\end{solution}

\begin{exercise}\label{b5:ex:05-s3-ring}
	设 \(S_3\) 的不可约特征标为 \(1,\varepsilon,\chi\)，次数为 \(1,1,2\)。计算其表示环乘法，并判断 \(\chi-\varepsilon\) 是否为实际特征标。
\end{exercise}
\begin{solution}
	在单位、换位、三循环三类上，\(\varepsilon=(1,-1,1)\)、\(\chi=(2,0,-1)\)。逐点乘法给出
	\[
		\varepsilon^2=1,\qquad\varepsilon\chi=\chi,\qquad \chi^2=1+\varepsilon+\chi.
	\]
	三者为加法基，故这些关系确定整个环。\(\chi-\varepsilon\) 次数为 \(1\)，但 \(\langle\varepsilon,\chi-\varepsilon\rangle=-1\)，不可约系数为负，因此不是实际特征标。
\end{solution}

\begin{exercise}\label{b5:ex:05-cyclotomic-integer}
	令 \(\zeta=e^{2\pi i/5}\)，\(a=\zeta+\zeta^{-1}\)。求 \(a\) 的首一整系数最小多项式，并说明代数整数不必是整数。
\end{exercise}
\begin{solution}
	由 \(1+\zeta+\zeta^2+\zeta^3+\zeta^4=0\)，有 \(a^2=\zeta^2+\zeta^{-2}+2=1-a\)，故 \(a^2+a-1=0\)。多项式 \(T^2+T-1\) 的判别式为 \(5\)，不是有理数平方，因此在 \(\mathbb Q[T]\) 中不可约。于是 \(a\) 是次数为 \(2\) 的代数整数而非有理数，当然也非整数。
\end{solution}

\begin{exercise}\label{b5:ex:05-class-sums-s3}
	设 \(D_8=\langle r,s:r^4=s^2=1,\ srs=r^{-1}\rangle\)，在 \(\mathbb Z[D_8]\) 中令 \(z=r^2\)、\(C=r+r^3\)、\(T=s+r^2s\)、\(U=rs+r^3s\)。证明它们为中心元素，计算 \(C^2,T^2,U^2,CT,CU,TU\)，并求三个类和在表示
	\[
		\rho(r)=\begin{pmatrix}0&-1\\1&0\end{pmatrix},\qquad
		\rho(s)=\begin{pmatrix}1&0\\0&-1\end{pmatrix}
	\]
	中的作用，核验所求标量满足这些整系数乘法关系。
\end{exercise}
\begin{solution}
	\(z\) 居中，\(\{r,r^3\}\)、\(\{s,r^2s\}\)、\(\{rs,r^3s\}\) 分别为共轭类，所以其类和居中。用 \(sr=r^{-1}s\) 展开乘积，得到
	\[
		\begin{aligned}
			C^2=T^2=U^2&=2(1+z),\\
			CT&=2U,\qquad CU=2T,\qquad TU=2C.
		\end{aligned}
	\]
	例如 \(C^2=r^2+2+r^6=2(1+z)\)，\(T^2=(1+z)^2=2(1+z)\)，而 \(T\cdot U=(1+z)^2srs=2(1+z)r^{-1}=2(r+r^3)\)。给定矩阵满足群的关系。\(\rho(r)\) 的两条复特征直线由 \(\rho(s)\) 交换，故表示不可约。又 \(\rho(z)=-I_2\)、\(\rho(r^3)=-\rho(r)\)，所以 \(\rho(C)=0\)；由 \(\rho(T)=(I_2+\rho(z))\rho(s)\)、\(\rho(U)=(I_2+\rho(z))\rho(rs)\)，另两个类和也作用为零。右边的 \(1+z\) 同样作用为零，故全部乘法关系在此表示中变为相应的标量等式。
\end{solution}

\begin{exercise}\label{b5:ex:05-p-group-degrees}
	设 \(|P|=p^a\)。证明每个复不可约表示的次数都是 \(p\) 的幂；若 \(P\ne1\) 且存在次数 \(d>1\) 的不可约表示，说明为什么 \(d^2<|P|\)。
\end{exercise}
\begin{solution}
	由\cref{b5:cor:character-degree-divisibility}，\(d\mid p^a\)，所以 \(d=p^b\)。由正则分解，全部不可约次数的平方和为 \(|P|\)。除了所取不可约表示，还至少有平凡表示，其平方贡献为 \(1\)，故 \(d^2+1\le|P|\)，从而严格小于 \(|P|\)。
\end{solution}

\begin{exercise}\label{b5:ex:05-rational-values-no-real-model}
	令 \(\theta\) 为 \(Q_8\) 的二维不可约复特征标。构造一个四维有理表示，其复化的特征标为 \(2\theta\)，并证明它在 \(\mathbb Q\) 上不可约。若一个非零有理表示的复化只含 \(\theta\) 这一不可约类型，求其最小可能维数。
\end{exercise}
\begin{solution}
	在基 \(1,i,j,k\) 下，用有理四元数代数的左乘给出矩阵
	\[
		A=\begin{pmatrix}0&-1&0&0\\1&0&0&0\\0&0&0&-1\\0&0&1&0\end{pmatrix},\qquad
		B=\begin{pmatrix}0&0&-1&0\\0&0&0&1\\1&0&0&0\\0&-1&0&0\end{pmatrix}.
	\]
	直接相乘可验 \(A^2=B^2=-I_4\)、\(AB=-BA\)，因此 \(i\mapsto A,j\mapsto B\) 满足 \(Q_8\) 的关系。其特征标在 \(1,-1\) 处为 \(4,-4\)，而 \(\pm A,\pm B,\pm AB\) 的迹为零，故复化的特征标为 \(2\theta\)。由复表示的完全可约性及特征标决定同构类，复化同构于两个 \(\theta\) 型表示的直和。

	若有非零真有理子表示，其复化是这个直和的非零真子表示，必只含一份 \(\theta\)，所以维数为 \(2\)。这与\cref{b5:prop:rational-values-no-real-model}排除二维有理模型的结论矛盾。故所构造的四维有理表示不可约。一般所求表示的复化特征标为 \(m\theta\)、\(m\ge1\)，维数为 \(2m\)；\(m=1\) 不可能，而本构造实现 \(m=2\)，故最小维数为 \(4\)。
\end{solution}

\begin{exercise}\label{b5:ex:05-klein-exact-index}
	令 \(G=C_2\times C_2\)，其四个不可约复特征标为 \(\lambda_1,\ldots,\lambda_4\)。对实际特征标 \(\psi=\sum_i n_i\lambda_i\)、\(n_i\ge0\)，证明：\(\psi\) 能写成循环子群的一维特征标的诱导的非负整数线性组合，当且仅当 \(\sum_i n_i\) 为偶数且 \(2n_i\le\sum_j n_j\) 对每个 \(i\) 成立。给出偶数次数但不满足这一要求的实际特征标。
\end{exercise}
\begin{solution}
	由\cref{b5:prop:klein-cyclic-induction-index}的诱导计算，每个 \(\lambda_i+\lambda_j\)、\(i\ne j\)，来自某个二阶子群的一维特征标。平凡子群的诱导为四个特征标之和，也能分为两个这样的配对。因此所求非负组合恰由这些两两配对生成。每个配对总次数为 \(2\)，且同一种类型至多占其中一份，故所列条件必要。

	反过来，写 \(\sum_i n_i=2s\)，按 \(s\) 归纳。\(s=0\) 时取空和。\(s>0\) 时至少两个 \(n_i\) 非零；将它们从大到小排列，取最大的两个，各减去 \(1\)，对应先取一个配对。余下总和为 \(2(s-1)\)。原来最大的两个都不超过 \(s\)，减一后不超过 \(s-1\)；第三大的数也不超过 \(s-1\)，否则前三个数至少有 \(3s>2s\) 的总和。故余下仍满足条件，由归纳可继续配对。最后，\(2\lambda_1\) 有偶数次数，却不能如此非负诱导表达：其第一种类型的重数 \(2\) 大于总次数的一半。它仍属于循环诱导的整数格；负系数与非负系数的要求不同。
\end{solution}

\begin{exercise}\label{b5:ex:05-elementary-abelian-denominator}
	设 \(G=C_p^r\)、\(r\ge2\)，\(\mathcal L\) 为其全部阶为 \(p\) 的子群。求使 \(N1_G\in I_{\mathrm{cyc}}(G)\) 成立的最小正整数 \(N\)，并给出一个表达。
\end{exercise}
\begin{solution}
	每个非平凡循环子群的指数为 \(p^{r-1}\)，平凡子群的指数为 \(p^r\)，所以所有循环诱导的次数都被 \(p^{r-1}\) 整除，必要条件是 \(p^{r-1}\mid N\)。子群数 \(L=|\mathcal L|=(p^r-1)/(p-1)\)，每个非单位元恰属一个这样的子群。于是逐点可验
	\[
		p^{r-1}1_G =\sum_{H\in\mathcal L}\operatorname{Ind}_H^G1_H -\frac{L-1}{p}\operatorname{Ind}_1^G1.
	\]
	在非单位元处两侧均为 \(p^{r-1}\)，在单位元处右侧为 \(Lp^{r-1}-(L-1)p^{r-1}\)。又 \((L-1)/p\) 为整数，故最小值正是 \(p^{r-1}\)。
\end{solution}

\begin{exercise}\label{b5:ex:05-permutation-artin-s3}
	只使用 \(S_3\) 的子群 \(1,C_2,C_3\) 的平凡特征标诱导，求平凡特征标 \(1\) 的有理表达。说明即使所有不可约特征标均为有理值，整数系数仍可能不足。
\end{exercise}
\begin{solution}
	记三个诱导分别为 \(\rho,\pi_2,\pi_3\)，则 \(\rho=1+\varepsilon+2\chi\)、\(\pi_2=1+\chi\)、\(\pi_3=1+\varepsilon\)。故
	\[
		1=\pi_2+\tfrac12\pi_3-\tfrac12\rho.
	\]
	若 \(1=a\rho+b\pi_2+c\pi_3\) 且 \(a,b,c\in\mathbb Z\)，比较 \(\chi,\varepsilon\) 系数得 \(b=-2a,c=-a\)，再比较平凡系数得 \(-2a=1\)，矛盾。允许非平凡子群特征标后，\eqref{b5:eq:s3-brauer-negative}才给出整数表达。
\end{solution}

\begin{exercise}\label{b5:ex:05-rational-classes-c5}
	求 \(C_5\) 上在有理共轭类中常值的有理函数空间，并用置换特征标写出任意这样的函数。比较它与全部复类函数空间的维数。
\end{exercise}
\begin{solution}
	有理共轭类只有 \(\{1\}\) 与四个非单位元组成的集合，故该有理空间维数为 \(2\)。若函数在这两处的值分别为 \(a,b\in\mathbb Q\)，则
	\[
		f=b\,1_{C_5}+\frac{a-b}{5}\operatorname{Ind}_1^{C_5}1.
	\]
	而 \(C_5\) 交换，每个元素单独成共轭类，全部复类函数空间维数为 \(5\)。减少到两个参数来自有理共轭条件，不是来自循环群特征标不足。
\end{solution}

\begin{exercise}\label{b5:ex:05-sylow-cover}
	分别对 \(S_3\) 与 \(C_6\)，判断只从全部 Sylow 子群的表示诱导，能否有理张成整个复表示环。
\end{exercise}
\begin{solution}
	\(S_3\) 的元素阶只有 \(1,2,3\)，每个元素都属于某个 Sylow 子群，且全部 Sylow 子群对共轭封闭。由\cref{b5:prop:subgroup-cover-induction}，有理张成成立。\(C_6\) 的生成元不属于二阶或三阶 Sylow 子群，且群交换，不能通过共轭进入它们。全部所述诱导在这个生成元处为零，故不能张成平凡特征标。
\end{solution}

\begin{exercise}\label{b5:ex:05-elementary-examples}
	对 \(C_6\)、\(C_2^2\times C_3\)、\(Q_8\)、\(S_3\)，分别判断它们是否为 \(2\)-初等群、\(3\)-初等群。
\end{exercise}
\begin{solution}
	\(C_6=C_2\times C_3\) 两种都是。\(C_2^2\times C_3\) 为 \(2\)-初等群，但不是 \(3\)-初等群，因为其与 \(3\) 互素的部分为非循环的 \(C_2^2\)。\(Q_8\) 为 \(2\)-初等群；若为 \(3\)-初等群，其 \(3\) 部分只能平凡，从而全群必须循环，矛盾。\(S_3\) 两种都不是：若有非平凡互素循环直因子，则该因子位于中心，而 \(Z(S_3)=1\)；若此因子平凡，全群又须为素数幂阶群。
\end{solution}

\begin{exercise}\label{b5:ex:05-elementary-quotient}
	证明有限 \(p\)-初等群的任意商群仍为 \(p\)-初等群；并说明同样的说法不能把“直积”随意改成“半直积”后仍当作定义。
\end{exercise}
\begin{solution}
	设 \(E=P\times C\)，\(N\triangleleft E\)。由\cref{b5:prop:elementary-subgroups}，\(N=(N\cap P)\times(N\cap C)\)，两个交分别在对应因子中正规，故
	\[
		E/N\simeq P/(N\cap P)\times C/(N\cap C).
	\]
	第一因子为 \(p\)-群，第二因子循环且阶与 \(p\) 互素。另一方面，\(S_3=C_3\rtimes C_2\) 是非平凡半直积，却不是初等群；交换两个因子的要求在定义及证明中都有实质作用。
\end{solution}

\begin{exercise}\label{b5:ex:05-p-regular-s3}
	在 \(S_3\) 中，对 \(p=2,3\) 分别列出各共轭类的 \(p\)-正则部分，并对三个不可约特征标验证模 \(p\) 约化公式。
\end{exercise}
\begin{solution}
	按 \(1,(12),(123)\) 排列，三个特征标为 \(1=(1,1,1)\)、\(\varepsilon=(1,-1,1)\)、\(\chi=(2,0,-1)\)，均为整数值。\(p=2\) 时换位的正则部分为 \(1\)，而三循环保持不变；于是需比较换位与单位元两列，逐项差为 \(0,-2,-2\)，均被 \(2\) 整除。\(p=3\) 时三循环的正则部分为 \(1\)，换位不变；三循环列与单位元列之差为 \(0,0,-3\)，均被 \(3\) 整除。
\end{solution}

\begin{exercise}\label{b5:ex:05-local-indices-klein}
	令 \(G=V\times C_3\)、\(V=C_2^2\)，计算 \(R(G)/I_p(G)\) 在 \(p=2\) 和奇素数 \(p\) 时的同构类型与指数，说明三个 \(C_3\) 特征标分别怎样影响循环诱导格。
\end{exercise}
\begin{solution}
	\(p=2\) 时 \(G\) 自身为 \(2\)-初等群，所以 \(I_2(G)=R(G)\)，商群为零，指数为 \(1\)。若 \(p=3\)，其 \(p\)-初等子群的互素循环部分只能是 \(V\) 的阶为 \(1\) 或 \(2\) 的子群；若 \(p>3\)，\(p\)-部分平凡，允许的子群就是循环子群。两种情形中，允许的子群都包含于某个 \(H\times C_3\)，其中 \(H\le V\) 为循环子群，而每个这样的直积本身循环，亦在允许范围内。

	令 \(\nu\) 为 \(C_3\) 的一个非平凡一维特征标。\(G\) 交换，全部不可约特征标为 \(\lambda\otimes\nu^j\)、\(\lambda\in\operatorname{Irr}(V)\)、\(j=0,1,2\)。由诱导的陪集模型，
	\[
		\operatorname{Ind}_{H\times C_3}^{V\times C_3}(\mu\otimes\nu^j) =(\operatorname{Ind}_H^V\mu)\otimes\nu^j.
	\]
	从更小的允许子群诱导时，先诱导到所含的 \(H\times C_3\)，用完全可约分解及诱导传递性，仍得到这些生成元的整数线性组合。因此在按 \(\nu^j\) 分开的三个加法直和项内，诱导格分别就是 \(I_{\mathrm{cyc}}(V)\)。由\cref{b5:prop:klein-cyclic-induction-index}，每一项的商为 \(\mathbb Z/2\mathbb Z\)，故奇素数情形
	\[
		R(G)/I_p(G)\cong(\mathbb Z/2\mathbb Z)^3, \qquad [R(G):I_p(G)]=8.
	\]
	三个 \(C_3\) 类型各留下一个独立的奇偶约束。这个指数与奇素数 \(p\) 互素，所以 \(p\)-局部化后商群消失，与逐素数生成结论相容。
\end{solution}

\begin{exercise}\label{b5:ex:05-fourier-choice}
	设 \(C=\langle c\rangle\) 阶为 \(6\)。在特征 \(5\) 的域 \(F\) 中取六阶单位根 \(z\)，令 \(\lambda_j(c)=z^j\)、\(0\le j<6\)，并取 \(T=\{1,c^2,c^4\}\)。计算 \(S_j=\sum_{t\in T}\lambda_j(t)\)，再用
	\[
		\frac16\sum_{j=0}^5S_j\lambda_j(c^a)^{-1}
	\]
	恢复 \(T\) 的示性函数。若改为特征 \(2\)，说明非空集合 \(T'=\{1,c^3\}\) 为何在所有一维 \(F\)-值特征标下都给出零和。
\end{exercise}
\begin{solution}
	\(S_j=1+z^{2j}+z^{4j}\)。若 \(3\mid j\)，该值为 \(3\)；否则 \(z^{2j}\) 为非平凡三阶单位根，等比求和给出 \(S_j=0\)。故只有 \(S_0=S_3=3\) 非零，恢复公式为
	\[
		\frac16(3+3z^{-3a})=\frac{1+(-1)^a}{2}.
	\]
	\(6\) 和 \(2\) 在 \(F\) 中都可逆，且 \(z^3=-1\)，所以此值在 \(a\) 为偶数时等于 \(1\)，为奇数时等于 \(0\)，正是 \(T\) 的示性函数。

	在特征 \(2\) 的域中，若 \(\lambda:C\rightarrow F^\times\) 为一维特征标，则 \(\lambda(c)^6=1\)。由 \(X^6-1=(X^3-1)^2\)，域中必有 \(\lambda(c)^3=1\)，所以
	\[
		\sum_{t\in T'}\lambda(t)=1+\lambda(c)^3=0.
	\]
	此时 \(|C|\) 在系数域中为零，不能使用含 \(1/|C|\) 的逆变换。非空集合可被所有一维特征标同时漏掉，说明 Brauer 证明要求所用循环群阶与剩余特征互素。
\end{solution}

\begin{exercise}\label{b5:ex:05-finiteness-lifting}
	用包含 \(\mathbb Z\subset\mathbb Q\) 说明：极大理想提升引理中的有限生成假设不能直接删去。这里是否还有元素满足的首一整方程？
\end{exercise}
\begin{solution}
	\(\mathbb Q\) 是域，唯一极大理想为 \(0\)，在 \(\mathbb Z\) 中的收缩仍为 \(0\)，并非 \((p)\)。因此任意素数理想 \((p)\subset\mathbb Z\) 都不能提升为 \(\mathbb Q\) 的极大理想。环 \(\mathbb Q\) 不是有限生成的 \(\mathbb Z\)-模，也不整于 \(\mathbb Z\)：例如 \(1/p\) 不是代数整数。问题不是模 \((p)\) 后选错了理想，而是 \(p\mathbb Q=\mathbb Q\)，此时整个纤维已经为零。
\end{solution}

\begin{exercise}\label{b5:ex:05-local-detection-primes}
	设 \(A\) 为有限交换群。证明 \(A_{(p)}\) 恰保留 \(A\) 的 \(p\)-主分量，并用一个例子说明只检测固定的一个素数不能推出 \(A=0\)。
\end{exercise}
\begin{solution}
	若元素阶与 \(p\) 互素，该阶在 \(\mathbb Z_{(p)}\) 中可逆，所以元素局部化后为零。若元素阶为 \(p^r\)，则与 \(p\) 互素的整数在这个循环群上乘法可逆，故不会使它消失。由有限交换群的主分量分解，\(A_{(p)}\) 与 \(p\)-主分量同构。也可直接对 \(|A|=p^r s\)、\(p\nmid s\)，用 Bézout 等式把每个元素拆为分别被 \(p^r\) 和 \(s\) 零化的两部分。若 \(q\ne p\)，非零群 \(\mathbb Z/q\mathbb Z\) 在 \(p\) 处局部化为零，说明须逐素数检测。
\end{solution}

\begin{exercise}\label{b5:ex:05-positive-restrictions}
	设 \(S_3\) 的不可约复特征标为 \(1,\varepsilon,\chi\)，其中后两者分别为符号特征标和二维特征标。对整数 \(a,b,c\)，给出虚特征标 \(f=a+b\varepsilon+c\chi\) 限制到每个初等子群均为实际特征标的充要条件。若 \(f\) 自身不是实际特征标，求满足这些条件的最小次数以及达到该次数的全部 \(f\)。
\end{exercise}
\begin{solution}
	\(S_3\) 的初等子群是平凡子群、二阶子群与三阶子群。在 \(C_2\) 上写非平凡一维特征标为 \(\eta\)，在 \(C_3\) 上写为 \(\lambda\)，则
	\[
		f|_{C_2}=(a+c)1+(b+c)\eta, \qquad f|_{C_3}=(a+b)1+c\lambda+c\lambda^2.
	\]
	实际特征标的不可约系数必须为非负整数，所以所求条件为
	\[
		a+c\ge0,\quad b+c\ge0,\quad a+b\ge0,\quad c\ge0.
	\]
	它们也保证平凡子群处的次数 \(a+b+2c\ge0\)，故无需另加条件。

	若全群上的 \(f\) 非实际，必须 \(a<0\) 或 \(b<0\)，且两者不能同时为负。设 \(a=-t\)、\(t\ge1\)，则 \(b\ge t\)、\(c\ge t\)，所以
	\[
		f(1)=-t+b+2c\ge2t\ge2.
	\]
	等号 \(f(1)=2\) 仅在 \(t=b=c=1\) 时成立。交换 \(a,b\) 的情形相同，故最小次数为 \(2\)，达到者恰为 \(-1+\varepsilon+\chi\) 和 \(1-\varepsilon+\chi\)。
\end{solution}

\begin{exercise}\label{b5:ex:05-quaternion-monomial}
	分类四元数群 \(Q_8\) 的二维不可约特征标 \(\theta\) 的全部一维诱导来源，即求出所有 \((H,\lambda)\)，使 \(H\le Q_8\)、\(\lambda\) 一维且 \(\operatorname{Ind}_H^{Q_8}\lambda=\theta\)。说明其他一维特征标在这些子群上诱导后如何分解。
\end{exercise}
\begin{solution}
	比较次数得 \([Q_8:H]=2\)，故 \(|H|=4\)。这样的子群恰为 \(\langle i\rangle,\langle j\rangle,\langle k\rangle\)，都循环。先取 \(H=\langle i\rangle\)，陪集代表为 \(1,j\)，且 \(jij^{-1}=i^{-1}\)。诱导特征标在 \(H\) 上为 \(\lambda+\lambda^{-1}\)，在其外为零。若 \(\lambda(i)=\pm\sqrt{-1}\)，则在 \(1,-1\) 处取值 \(2,-2\)，在其余处为零，恰为 \(\theta\)。两个另外的子群同理，各有两个忠实一维特征标可用，故共有六组来源。

	若 \(\lambda(i)=1\) 或 \(-1\)，它不是忠实特征标，并且共轭固定它。它在 \(Q_8\) 上恰有两个一维延拓：令 \(i\) 的像为所给的 \(\pm1\)，令 \(j\) 的像分别为 \(1,-1\)，这些赋值均满足群的关系。由 Frobenius 互反，两种延拓各出现一次，且次数和已经为 \(2\)，所以诱导恰为这两个一维特征标的直和，不再不可约。对 \(\langle j\rangle,\langle k\rangle\) 的结论由相同计算得到。
\end{solution}

\begin{exercise}\label{b5:ex:05-product-monomial}
	令 \(E=Q_8\times C_3\)，\(\theta\) 为 \(Q_8\) 的二维不可约特征标，\(\nu\) 为 \(C_3\) 的非平凡一维特征标。写出 \(\theta\otimes\nu\) 的一维诱导表达及全部特征标值。
\end{exercise}
\begin{solution}
	取 \(H=\langle i\rangle\)、\(\lambda(i)=\sqrt{-1}\)，由陪集基或\cref{b5:prop:elementary-monomial}的构造，
	\[
		\theta\otimes\nu =\operatorname{Ind}_{H\times C_3}^{Q_8\times C_3}(\lambda\otimes\nu).
	\]
	对子群 \(H\times C_3\)，\(H\) 为 \(2\)-群且 \(C_3\) 为互素循环因子，仍是允许的 \(2\)-初等群。若 \(c\) 为 \(C_3\) 生成元、\(\nu(c)=\zeta_3\)，则特征标在 \((1,c^j)\) 处为 \(2\zeta_3^j\)，在 \((-1,c^j)\) 处为 \(-2\zeta_3^j\)，在其余元素处为零。
\end{solution}

\begin{exercise}\label{b5:ex:05-sign-negative-induction}
	用 \(S_3\) 的二阶、三阶子群的一维特征标整诱导表达符号特征标 \(\varepsilon\)，并证明任何这样的表达都不可能只含非负系数。
\end{exercise}
\begin{solution}
	令 \(\eta\) 为 \(C_2\) 的非平凡一维特征标，\(\lambda\) 为 \(C_3\) 的非平凡一维特征标。Frobenius 互反给出 \(\operatorname{Ind}_{C_2}^{S_3}\eta=\varepsilon+\chi\)，而 \(\operatorname{Ind}_{C_3}^{S_3}\lambda=\chi\)。故
	\[
		\varepsilon=\operatorname{Ind}_{C_2}^{S_3}\eta -\operatorname{Ind}_{C_3}^{S_3}\lambda.
	\]
	\(S_3\) 的初等子群都是真子群，一维诱导的次数至少为 \(2\)，所以任何非零非负整数线性组合的次数至少为 \(2\)，不能等于 \(\varepsilon(1)=1\)。
\end{solution}

\begin{exercise}\label{b5:ex:05-nonintegral-average}
	证明两个单位根的平均 \((1+i)/2\) 不是代数整数，并计算它的最小多项式及共轭乘积。
\end{exercise}
\begin{solution}
	设 \(z=(1+i)/2\)，则 \(z^2-z+1/2=0\)。判别式为 \(-1\)，该首一二次多项式在 \(\mathbb Q[T]\) 中不可约，所以是最小多项式。其常数项为 \(1/2\)，不是整数；由\cref{b5:lem:algebraic-integer-arithmetic}，\(z\) 不可能是代数整数。两个共轭根的乘积正是 \((1+i)(1-i)/4=1/2\)，展示了平均值证明中整数乘积条件不可缺少。
\end{solution}

\begin{exercise}\label{b5:ex:05-a4-coprime-failure}
	从 \(A_4\) 在四个字母上的置换表示中除去平凡表示，得次数为 \(3\) 的特征标 \(\chi\)。证明它不可约；再用双换位说明消失或标量作用引理中的互素假设不能删除。
\end{exercise}
\begin{solution}
	在单位、三个双换位、八个三循环处，\(\chi\) 的值分别为 \(3,-1,0\)。于是 \(\langle\chi,\chi\rangle=(9+3)/12=1\)，故不可约。双换位在四维置换表示上的特征值为 \(1,1,-1,-1\)，去掉一个平凡方向后为 \(1,-1,-1\)。所以其迹为 \(-1\ne0\)，矩阵亦非标量。双换位的共轭类大小为 \(3\)，与次数 \(3\) 不互素，正落在引理假设之外。
\end{solution}

\begin{exercise}\label{b5:ex:05-dihedral-class-normal}
	设 \(D_{14}=\langle r,s:r^7=s^2=1,\ srs=r^{-1}\rangle\)，下标表示群阶。求反射的共轭类，并计算由这个类中元素之商生成的正规子群。
\end{exercise}
\begin{solution}
	由 \(r^a s r^{-a}=r^{2a}s\)，且 \(2\) 在模 \(7\) 下可逆，七个反射 \(r^j s\) 构成同一个共轭类 \(K\)，大小为 \(7\)。两个反射之商为 \((r^i s)(r^j s)^{-1}=r^{i-j}\)，遍历全部旋转。因此 \(\langle xy^{-1}:x,y\in K\rangle=\langle r\rangle\)，是阶为 \(7\) 的非平凡真正规子群。把旋转送到 \(1\)、反射送到 \(-1\) 的一维表示正好零化这个子群，与正文中的标量作用构造一致。
\end{solution}

\begin{exercise}\label{b5:ex:05-perfect-two-primes}
	设有限群 \(G\) 满足 \(G'=G\)，且 \(|G|\) 至多含两个不同素因子。证明 \(G=1\)。再证明每个非平凡有限可解群都有非平凡一维复特征标。
\end{exercise}
\begin{solution}
	由\cref{b5:thm:burnside-solvability}，\(G\) 可解。但 \(G'=G\) 使全部导出项都等于 \(G\)，只有 \(G=1\) 才能终止于 \(1\)。对任意非平凡有限可解群，导出群必须是真子群，所以交换化 \(G/G'\) 非平凡；取其一个素数阶商并选非平凡一维复特征标，提升到 \(G\) 就得到非平凡一维特征标。因此非平凡有限可解群并非只具有次数大于一的非平凡特征标。
\end{solution}

\begin{exercise}\label{b5:ex:05-a4-derived}
	求 \(A_4\) 的导出列，并说明 \(p^a q^b\) 阶群的可解性不保证全部 Sylow 子群正规。
\end{exercise}
\begin{solution}
	Klein 四元子群 \(V_4\) 由单位与三个双换位组成，正规且商 \(A_4/V_4\simeq C_3\)，故 \(A_4'\subset V_4\)。\(A_4\) 不交换，所以 \(A_4'\ne1\)；三个双换位在 \(A_4\) 中共轭，\(V_4\) 内任何包含其中一个的正规子群必须包含全部三个。因此 \(A_4'=V_4\)，再有 \(V_4'=1\)。导出列为 \(A_4,V_4,1\)。八个三循环组成四个不同三阶子群，故 Sylow \(3\)-子群不正规，尽管 \(|A_4|=2^2\cdot3\) 且群可解。
\end{solution}

\begin{exercise}\label{b5:ex:05-minimal-normal-linear}
	设 \(p\) 为素数、\(r\ge1\) 为整数，\(N=(\mathbb F_p^r,+)\)、\(H\le\operatorname{GL}_r(\mathbb F_p)\)。令 \(G=N\rtimes H\)，乘法为 \((v,h)(w,k)=(v+hw,hk)\)。证明 \(N\) 为 \(G\) 的极小非平凡正规子群，当且仅当 \(H\) 在 \(\mathbb F_p^r\) 上的作用不可约，并求 \(C_G(N)\)。对 \(p=3,r=2\) 及
	\[
		H=\left\langle\begin{pmatrix}0&-1\\1&0\end{pmatrix}\right\rangle
	\]
	完成具体计算，说明得到的极小正规子群为什么不是简单群。
\end{exercise}
\begin{solution}
	\(N\) 是投影 \(G\rightarrow H\) 的核，所以正规，且
	\[
		(v,h)(w,1)(v,h)^{-1}=(hw,1).
	\]
	\(N\) 的子群恰是 \(\mathbb F_p\)-子空间，因而其中在 \(G\) 中正规的子群恰是 \(H\)-不变子空间。这给出极小正规性与不可约性的等价。所示共轭公式也说明，\((v,h)\) 中心化全部 \(N\) 当且仅当 \(h\) 为恒等矩阵；\(H\) 已作为矩阵子群嵌入，故 \(C_G(N)=N\)。因此 \(G\) 对 \(N\) 的共轭作用的核为 \(N\)，商群 \(G/N\cong H\) 的作用忠实。

	具体矩阵 \(A\) 在 \(\mathbb F_3\) 上满足 \(A^2=-I\)、\(A^4=I\)，所以 \(H\cong C_4\)。多项式 \(X^2+1\) 在 \(\mathbb F_3\) 中无根，故 \(A\) 没有特征直线，二维作用不可约。因此阶为 \(9\) 的 \(N\cong C_3^2\) 是阶为 \(36\) 的 \(G\) 的极小非平凡正规子群。但 \(N\) 自身交换，任意一维子空间都是其非平凡真正规子群，所以 \(N\) 不是简单群。这些直线都不被整个 \(H\) 保持，也就不在 \(G\) 中正规。
\end{solution}

\begin{exercise}\label{b5:ex:05-real-quaternion-density}
	令实四元数除环 \(\mathbb H\) 通过左乘作用于实向量空间 \(\mathbb H\)。证明这是不可约实矩阵子代数，但不等于 \(\operatorname{End}_{\mathbb R}(\mathbb H)\)。
\end{exercise}
\begin{solution}
	若非零实子空间 \(W\) 被全部左乘保持，取 \(0\ne w\in W\)。对任意 \(y\in\mathbb H\)，因四元数可逆，\(y=(yw^{-1})w\in W\)，所以 \(W=\mathbb H\)。故左乘代数作用不可约。该子代数的实维数为 \(4\)，而全部自同态代数实维数为 \(4^2=16\)，两者不相等。这是底域未代数闭时的失败情形；除环的非标量右乘也解释了标量交换子结论为何不成立。
\end{solution}

\begin{exercise}\label{b5:ex:05-explicit-matrix-generation}
	设 \(k\) 为含 \(n\) 个不同元素 \(a_1,\ldots,a_n\) 的域。令 \(D=\operatorname{diag}(a_1,\ldots,a_n)\)，令 \(S\) 循环置换标准基。直接证明由 \(D,S\) 生成的含幺 \(k\)-代数为 \(M_n(k)\)。
\end{exercise}
\begin{solution}
	Lagrange 插值多项式\footnote{拉格朗日（Giuseppe Lodovico Lagrangia，后用 Joseph-Louis Lagrange，1736-1813），意大利裔法国数学家，发展变分法，并著有《分析力学》。}
	\[
		f_i(T)=\prod_{j\ne i}\frac{T-a_j}{a_i-a_j}
	\]
	满足 \(f_i(D)=E_{ii}\)，所以所有坐标投影均在代数中。对任意 \(i,j\)，选 \(0\le r<n\) 使 \(S^r e_j=e_i\)，则 \(E_{ii}S^rE_{jj}=E_{ij}\)。全部矩阵单位都在所生成代数中，而它们是 \(M_n(k)\) 的基，故结论成立。
\end{solution}

\begin{exercise}\label{b5:ex:05-necklace-count}
	设 \(\ell\) 为素数，\(C_\ell\) 通过循环旋转作用在长为 \(\ell\) 的二进制串上。求轨道数，并解释它是某个置换表示中平凡表示的重数。
\end{exercise}
\begin{solution}
	单位旋转固定全部 \(2^\ell\) 个串。每个非单位旋转在位置集合上是一个 \(\ell\)-循环，所以固定串必须所有位相同，恰有两种。因此轨道数为
	\[
		\frac{2^\ell+2(\ell-1)}{\ell}.
	\]
	置换特征标在每个旋转处等于其固定串数，和 \(1_{C_\ell}\) 的内积便是上述平均，即平凡表示的重数。
\end{solution}

\begin{exercise}\label{b5:ex:05-scalar-subgroup}
	设 \(\rho:G\rightarrow\operatorname{GL}(V)\) 是有限群的非零复不可约表示，令 \(Z_\rho=\bigl\{g\in G:\rho(g)\;\text{为标量}\bigr\}\)。证明 \(Z_\rho\) 为正规子群，且 \(Z_\rho/\ker\rho=Z\bigl(G/\ker\rho\bigr)\)。忠实情形有何结论？
\end{exercise}
\begin{solution}
	非零标量矩阵组成 \(\operatorname{GL}(V)\) 的中心子群，其原像 \(Z_\rho\) 因而正规。若 \(\rho(g)\) 为标量，它与每个 \(\rho(h)\) 交换，故 \(g\ker\rho\) 位于商群中心。反过来，若这个陪集居中，则 \(\rho(g)\) 与全部表示矩阵交换，由 Schur 引理是标量。两边相等，且 \(\ker\rho\subset Z_\rho\)，故得到所述商群等式。若表示忠实，则 \(Z_\rho=Z(G)\)；在非交换单群的非平凡不可约表示中，标量元素只能是单位。
\end{solution}

\begin{exercise}\label{b5:ex:05-simple-character-constraint}
	设 \(G\) 为有限非交换单群，\(g\ne1\)。证明存在非平凡不可约复特征标 \(\chi\)，使 \(\gcd\bigl(\chi(1),|g^G|\bigr)>1\)。
\end{exercise}
\begin{solution}
	反设全部非平凡不可约次数均与 \(|g^G|\) 互素。每个非平凡不可约表示的核为正规真子群，故由单性为 \(1\)。由\cref{b5:lem:coprime-class-scalar}，若 \(\chi(g)\ne0\)，则 \(g\) 在一个忠实不可约表示中作用为标量，从而位于 \(Z(G)=1\)，矛盾。因此所有非平凡不可约特征标在 \(g\) 处均为零。正则特征标分解遂在 \(g\ne1\) 处给出 \(0=1\)，矛盾。这一结论约束类大小与次数的公共素因子，而不只是分别要求它们整除群阶。
\end{solution}
